	\begin{center}
		\section{计算几何}
	\end{center}

	\subsection{点线基础模板}

		以下是计算几何最底层的几个基础类型：符号判定 \texttt{sgn}、点 / 向量 \texttt{Point}（点和向量同一类型，向量当成「以原点为起点」的点）、叉乘 / 点乘、两点式直线 \texttt{Line}。\textbf{后续所有子节的代码都建立在这些定义之上}，每节顶部会以注释形式列出本节用到的依赖（具体实现见此处）。

		\texttt{Point} / \texttt{Line} 都是 \texttt{template<class T>} 泛型（无默认），调用方在 call site 显式给 \texttt{T}（如 \texttt{Point<ll>}、\texttt{Point<DB>}）；同一份模板可同时被整数 \texttt{Point<ll>} 与浮点 \texttt{Point<DB>} 实例化（对拍 generator 常见需求）。\textbf{向量类型不另起别名}，下游函数体里取边向量直接 \texttt{auto v = poly[i+1] - poly[i]}，让编译器把它推成 \texttt{Point<T>}（点和向量同形）。当 \texttt{T} 为\textbf{整数类型}（含 \texttt{ll} 与 \texttt{\_\_int128}）时，\texttt{cross} / \texttt{norm} 自动放大到 \texttt{\_\_int128\_t} 计算以避免溢出。\texttt{Point::abs()} 走自己写的 \texttt{sqrtL}：浮点直接 \texttt{sqrtl}，\textbf{整数走 \texttt{isqrt} + 修正}，避免 \texttt{(long double)\,i128} 在 $n > 2^{53}$ 时丢整数精度。

		\href{https://gist.github.com/znzryb/11d094702ccdb989efb8aed573ab0885}{change log}

		\begin{minted}{cpp}
// floor(sqrt(n))，n >= 0；二分 + 防溢出（mid <= n/mid 等价 mid*mid <= n）
i128 isqrt(i128 n) {
	if (n <= 0) return 0;
	i128 lo = 0, hi = (i128) 2e18; // sqrt(4e36) <= 2e18
	while (lo < hi) {
		i128 mid = lo + (hi - lo + 1) / 2;
		if (mid <= n / mid) lo = mid;
		else                hi = mid - 1;
	}
	return lo;
}

// 统一 sqrt 入口：浮点直接 sqrtl；整数（含 i128）走 isqrt + 修正，
// 避免 (long double) i128 在 n > 2^53 时丢整数精度
template<typename T>
long double sqrtL(T n) {
	if constexpr (is_floating_point_v<T>) {
		return sqrtl((long double) n);
	} else {
		i128 ni = (i128) n;
		if (ni <= 0) return 0;
		i128 i = isqrt(ni);
		long double id = (long double) i;
		long double rd = (long double) (ni - i * i);
		// sqrt(i^2 + r) = i * sqrt(1 + r/i^2)
		return id * sqrtl(1.0L + rd / (id * id));
	}
}

// x*x 安全平方: 整数升 i128 防溢出, 浮点原样 T*T (i128 没浮点意义).
template<class T>
auto squared(T x) {
	if constexpr (is_floating_point_v<T>) {
		return x * x;
	} else {
		return (i128) x * (i128) x;
	}
}

template<class T>
int sgn(T x) {
	if (-eps < x && x < eps) return 0;
	return x < 0 ? -1 : 1;
}

template<class T>
struct Point {
	T x, y;
	int id = -1; // 原始下标. -1 表示派生点 (+ - * 等算术结果) 或未赋值, id 无意义

	explicit Point(T x_ = 0, T y_ = 0) : x(x_), y(y_) {}

	// Point<U> -> Point<T> 隐式 (common_type 单调升档方向, U -> T 必须无损):
	// common_type_t<U, T> == T 才放行 (例 ll -> DB 通; DB -> ll 拦, 防 truncation)
	template<class U, class = enable_if_t<is_same_v<common_type_t<U, T>, T>>>
	Point(const Point<U> &o) : x((T) o.x), y((T) o.y), id(o.id) {}

	// 只对左值生效的复合运算符
	Point &operator+=(Point p) & { x += p.x, y += p.y; return *this; }
	Point &operator-=(Point p) & { x -= p.x, y -= p.y; return *this; }
	Point &operator*=(T t) & { x *= t, y *= t; return *this; }
	Point &operator/=(T t) & { x /= t, y /= t; return *this; }

	Point operator-() const { return Point(-x, -y); }

	friend Point operator+(Point a, Point b) { return a += b; }
	friend Point operator-(Point a, Point b) { return a -= b; }

	// 标量乘除: 返回 Point<common_type_t<T, U>> (C++14), 自动算公共类型
	// <ll,ll>=ll / <ll,DB>=DB / <DB,ll>=DB; R==T 走 copy ctor, R!=T 走 converting ctor
	template<class U>
	friend auto operator*(Point a, U b) {
		Point<common_type_t<T, U>> r = a; r *= b; return r;
	}
	template<class U>
	friend auto operator*(U a, Point b) {
		Point<common_type_t<T, U>> r = b; r *= a; return r;
	}
	template<class U>
	friend auto operator/(Point a, U b) {
		Point<common_type_t<T, U>> r = a; r /= b; return r;
	}

	friend bool operator<(Point a, Point b) {
		int sx = sgn(a.x - b.x);
		if (sx != 0) return a.x < b.x;
		return a.y < b.y;
	}
	friend bool operator>(Point a, Point b) { return b < a; }
	friend bool operator==(Point a, Point b) {
		return sgn(a.x - b.x) == 0 && sgn(a.y - b.y) == 0;
	}
	friend bool operator!=(Point a, Point b) { return !(a == b); }
	friend istream &operator>>(istream &is, Point &p) { return is >> p.x >> p.y; }

	// !is_floating_point_v 而不是 is_integral_v：让 i128 / __int128 也走整数分支
	auto norm() const {
		if constexpr (!is_floating_point_v<T>) {
			return (__int128_t) x * x + (__int128_t) y * y;
		} else {
			return x * x + y * y;
		}
	}
	DB abs() const { return sqrtL(norm()); }

	Point rot90() const { return Point(-y, x); }
	Point<DB> rot(DB ang) const {
		return Point<DB>(cosl(ang) * x - sinl(ang) * y,
		                 sinl(ang) * x + cosl(ang) * y);
	}

	// L1 范数 (曼哈顿) = |x| + |y|. 整型分支走 __int128_t 防溢出.
	auto normL1() const {
		auto ax = (x < 0 ? -x : x), ay = (y < 0 ? -y : y);
		if constexpr (!is_floating_point_v<T>) {
			return (__int128_t) ax + (__int128_t) ay;
		} else {
			return ax + ay;
		}
	}

	// Linf 范数 (切比雪夫) = max(|x|, |y|). 返回类型与 T 同.
	T normLinf() const {
		auto ax = (x < 0 ? -x : x), ay = (y < 0 ? -y : y);
		return ax > ay ? ax : ay;
	}
};

template<class T>
auto cross(Point<T> a, Point<T> b) {
	if constexpr (!is_floating_point_v<T>) {
		return (__int128_t) a.x * b.y - (__int128_t) a.y * b.x;
	} else {
		return a.x * b.y - a.y * b.x;
	}
}
// 三参数重载：以 o 为原点算 (a - o) x (b - o)
// 用途：判 3 点 turn (>0 左转 / <0 右转 / =0 共线) / 2 倍三角形有向面积
template<class T>
auto cross(Point<T> o, Point<T> a, Point<T> b) { return cross(a - o, b - o); }

template<class T, class U>
auto dot(Point<T> a, Point<U> b) {
	// 在 T 为整数的时候，使用 __int128_t 整型进行计算
	if constexpr (!is_floating_point_v<T> && !is_floating_point_v<U>) {
		return (__int128_t) a.x * b.x + (__int128_t) a.y * b.y;
	} else {
		return a.x * b.x + a.y * b.y;
	}
}

// 两点式直线
template<class T>
struct Line {
	Point<T> p1, p2;

	Line() = default;

	Line(Point<T> a, Point<T> b) : p1(a), p2(b) {}

	// Line<U> -> Line<T> 隐式 (common_type 单调升档方向, 同 Point converting ctor)
	template<class U, class = enable_if_t<is_same_v<common_type_t<U, T>, T>>>
	Line(const Line<U> &o) : p1(o.p1), p2(o.p2) {}

	// AC https://qoj.ac/submission/2146870
	Line(T k, T c) : p1(Point<T>(0, c)), p2(Point<T>(1, k + c)) {}

	friend bool operator<(Line a, Line b) {
		if (a.p1 != b.p1) return a.p1 < b.p1;
		return a.p2 < b.p2;
	}
	friend bool operator==(Line a, Line b) {
		return a.p1 == b.p1 && a.p2 == b.p2;
	}
	friend istream &operator>>(istream &is, Line &e) {
		T a, b, c, d;
		is >> a >> b >> c >> d;
		e = Line(Point<T>(a, b), Point<T>(c, d));
		return is;
	}
};

Line<DB> gen_line_from_general(DB a, DB b, DB c) {
	// 方向向量 (b, -a)
	Point<DB> p1, p2;
	if (b != 0) {
		p1 = Point<DB>(0, -c / b);
	} else {
		p1 = Point<DB>(-c / a, 0);
	}
	p2 = Point<DB>(p1.x + b, p1.y - a);
	return Line<DB>(p1, p2);
}

mt19937_64 rng(std::chrono::steady_clock::now().time_since_epoch().count());
		\end{minted}

		\textbf{参考题单}：\href{https://vjudge.net/article/12247}{计算几何模板题目总结}（vjudge）以及 \href{https://qoj.ac/contest/2646}{QOJ 模板赛 2646}。

		\subsection{点在线上的投影点}

		\href{https://vjudge.net/solution/64166651}{vjudge 提交}。

		设 \texttt{vec = l.p2 - l.p1}，把向量 \texttt{p - l.p1} 在 \texttt{vec} 方向上的归一化投影长度记为 $r$，于是投影点等于 \texttt{l.p1 + r * vec}。\textbf{返回值固定为 \texttt{Point<DB>}}（除法走浮点，输入 \texttt{T / U} 是 \texttt{ll} 还是 \texttt{DB} 都向 \texttt{DB} 收敛），不再 \texttt{assert} 浮点。

		\begin{minted}{cpp}
// 依赖（基础模板已定义）：
//   struct Point, struct Line
//   dot(Point, Point)
// 使用这个东西，Point 必须是 double 等浮点类型
template<class T, class U>
Point<DB> project(Point<T> p, Line<U> l) {
	Point<DB> vec = l.p2 - l.p1;
	//                        这个距离反正早除晚除都要除掉
	DB r = dot(vec, p - l.p1) / ((DB) vec.x * vec.x + vec.y * vec.y);
	return l.p1 + vec * r;
}
		\end{minted}

		\subsection{对称点}

		\href{https://vjudge.net/solution/64168723}{vjudge 提交}。

		设 \texttt{proj} 为投影点，则对称点 \texttt{= proj + proj - p}。直观上：\textcolor{blue}{\textbf{以投影点为中点}}，原点和对称点处于直线两侧、距离相等。

		\begin{center}
			\begin{tikzpicture}[>=Stealth]
				\draw[thick, blue!70!black] (-2.2, 0) -- (2.6, 0) node[right, font=\scriptsize] {$L$};
				\coordinate (P)  at (0.5,  1.4);
				\coordinate (Pp) at (0.5,  0);
				\coordinate (Pr) at (0.5, -1.4);
				\draw[dashed, gray] (P) -- (Pr);
				\draw[gray] (0.3, 0) -- (0.3, 0.2) -- (0.5, 0.2);
				\fill[red]    (P)  circle (1.8pt) node[above right, font=\scriptsize, color=red]    {$P$};
				\fill[teal]   (Pp) circle (1.8pt) node[above left,  font=\scriptsize, color=teal]   {$P'$ (proj)};
				\fill[orange] (Pr) circle (1.8pt) node[below right, font=\scriptsize, color=orange] {$P''$};
				\draw[<->, blue!50!black, thin] (-1.5, 0.05) -- (-1.5, 1.35) node[midway, left, font=\tiny] {$d$};
				\draw[<->, blue!50!black, thin] (-1.5, -0.05) -- (-1.5, -1.35) node[midway, left, font=\tiny] {$d$};
			\end{tikzpicture}

			\vspace{2pt}
			{\footnotesize\sffamily\color{gray} 图注：$P'$（投影）是 $P$ 到 $L$ 的垂足；$P''$（对称点）满足 $P'$ 是 $P, P''$ 中点，从而 $P'' = 2P' - P$。}
		\end{center}

		\begin{minted}{cpp}
// 依赖：
//   struct Point, struct Line
//   project(Point, Line)（见前一节）
template<class T, class U>
Point<DB> reflect(Point<T> p, Line<U> l) {
	Point<DB> proj = project(p, l);
	return proj + proj - p;
}
		\end{minted}

		\subsection{点线位置判断（CCW）}

		\href{https://vjudge.net/solution/64174955}{vjudge 提交}。

		给一条有向直线 $\overrightarrow{a \to b}$，目标点 $p$ 落在何处由 $\mathrm{cross}(\vec{ab}, \vec{ap})$ 与 $\mathrm{dot}(\vec{ab}, \vec{ap})$ 决定，共五种情况：

		\begin{center}
			\begin{tikzpicture}[>=Stealth, x=1cm, y=1cm]
				\begin{scope}
					\draw[->, thick] (0,0) -- (1.2,0);
					\fill (0,0) circle (1pt); \fill (1.2,0) circle (1pt);
					\fill[teal] (0.6, 0.5) circle (1.6pt) node[above, font=\tiny, color=teal] {$p$};
					\node[font=\tiny, anchor=north] at (0.6, -0.55) {COUNTER\_CCW};
				\end{scope}
				\begin{scope}[xshift=2.4cm]
					\draw[->, thick] (0,0) -- (1.2,0);
					\fill (0,0) circle (1pt); \fill (1.2,0) circle (1pt);
					\fill[orange] (0.6, -0.5) circle (1.6pt) node[below, font=\tiny, color=orange] {$p$};
					\node[font=\tiny, anchor=north] at (0.6, -0.85) {CLOCKWISE};
				\end{scope}
				\begin{scope}[xshift=4.8cm]
					\draw[->, thick] (0,0) -- (1.2,0);
					\fill (0,0) circle (1pt); \fill (1.2,0) circle (1pt);
					\fill[violet] (-0.5, 0) circle (1.6pt) node[above, font=\tiny, color=violet] {$p$};
					\node[font=\tiny, anchor=north] at (0.35, -0.55) {ONLINE\_BACK};
				\end{scope}

				\begin{scope}[yshift=-2.0cm]
					\draw[->, thick] (0,0) -- (1.2,0);
					\fill (0,0) circle (1pt); \fill (1.2,0) circle (1pt);
					\fill[red] (1.7, 0) circle (1.6pt) node[above, font=\tiny, color=red] {$p$};
					\node[font=\tiny, anchor=north] at (0.85, -0.55) {ONLINE\_FRONT};
				\end{scope}
				\begin{scope}[xshift=2.4cm, yshift=-2.0cm]
					\draw[->, thick] (0,0) -- (1.2,0);
					\fill (0,0) circle (1pt); \fill (1.2,0) circle (1pt);
					\fill[blue] (0.6, 0) circle (1.6pt) node[above, font=\tiny, color=blue] {$p$};
					\node[font=\tiny, anchor=north] at (0.6, -0.55) {ON\_SEGMENT};
				\end{scope}
				\node[font=\tiny] at (0, 0.18) {$a$};
				\node[font=\tiny] at (1.2, 0.18) {$b$};
			\end{tikzpicture}

			\vspace{2pt}
			{\footnotesize\sffamily\color{gray} 图注：先看 $\mathrm{cross}(\vec{ab},\vec{ap})$ 的符号——正左 / 负右；为零再看 $\mathrm{dot}$ 与 $|\vec{ab}|, |\vec{ap}|$ 的比较，区分 $p$ 落在线段反方向延长（BACK）、正方向延长（FRONT）还是线段内部（ON）。}
		\end{center}

		\begin{minted}{cpp}
// 依赖：
//   struct Point, struct Line
//   sgn, cross, dot
//   Point::norm()
constexpr int ON_SEGMENT        =  0;
constexpr int COUNTER_CLOCKWISE =  1;
constexpr int CLOCKWISE         = -1;
constexpr int ONLINE_BACK       =  2;
constexpr int ONLINE_FRONT      = -2;

template<class T, class U>
int CCW(Point<T> p, Line<U> l) {
	Point<DB> a = l.p2 - l.p1;
	Point<DB> b = p - l.p1;
	if (sgn(cross(a, b)) > 0) {
		return COUNTER_CLOCKWISE;
	}
	if (sgn(cross(a, b)) < 0) {
		return CLOCKWISE;
	}
	if (sgn(dot(a, b)) < 0) {
		return ONLINE_BACK;
	}
	if (a.norm() < b.norm()) return ONLINE_FRONT;
	return ON_SEGMENT;
}
		\end{minted}

		\subsubsection{逆时针弧角 \texttt{ccw\_angle}（标量 / 向量两版）}

		求「从 \texttt{from} 逆时针转到 \texttt{to}」的有向夹角，结果恒落在 $[0, 2\pi)$，而不是 $|\texttt{to} - \texttt{from}|$（后者在跨 $0/2\pi$ 处会算错）。两个重载：标量版输入两个极角，向量版输入两个方向向量、用单次 $\operatorname{atan2}(\mathrm{cross}, \mathrm{dot})$（叉积当 $\sin$ 分量、点积当 $\cos$ 分量）一步求出，比两次 \texttt{atan2} 相减更稳。

		\begin{minted}{cpp}
// 依赖：
//   struct Point, cross, dot（向量版用到）
// 圆弧 from -> to 的逆时针弧角差 [0, 2pi), 不是 abs(to - from)
static DB ccw_angle(DB theta_from, DB theta_to) {
	DB d = std::fmod(theta_to - theta_from, 2 * PI);
	if (d < 0) d += 2 * PI;
	return d;
}

// 向量版: 从方向向量 from 逆时针转到方向向量 to 的角 [0, 2pi).
// 用单次 atan2(cross, dot) (cross 当 sin 分量 / dot 当 cos 分量), 比两次 atan2 相减稳.
template<class T, class U>
static DB ccw_angle(Point<T> from, Point<U> to) {
	DB d = std::atan2((DB)cross(from, to), (DB)dot(from, to));
	if (d < 0) d += 2 * PI;
	return d;
}
		\end{minted}

		\subsection{判断两直线是否垂直 / 平行}

		\href{https://vjudge.net/solution/64175524}{vjudge 提交}。

		分别取两条直线的方向向量，看 \texttt{dot} 是否为零（垂直）或 \texttt{cross} 是否为零（平行）。

		\begin{minted}{cpp}
// 依赖：
//   struct Line
//   sgn, cross, dot
template<class T, class U>
bool isOrthogonal(Line<T> a, Line<U> b) {
	auto c = a.p2 - a.p1, d = b.p2 - b.p1;
	return !sgn(dot(c, d));
}
template<class T, class U>
bool isParallel(Line<T> a, Line<U> b) {
	auto c = a.p2 - a.p1, d = b.p2 - b.p1;
	return !sgn(cross(c, d));
}
		\end{minted}

		\subsection{判断两线段是否相交}

		\href{https://vjudge.net/solution/64175810}{vjudge 提交}。

		\texttt{a} 与 \texttt{b} 相交当且仅当 \texttt{b} 的两端点关于 \texttt{a} 异侧（含端点重合）\textbf{且} \texttt{a} 的两端点关于 \texttt{b} 异侧。利用 \texttt{CCW} 的符号相乘 $\le 0$ 即可。

		\begin{minted}{cpp}
// 依赖：
//   struct Line
//   CCW(Point, Line)
template<class T, class U>
bool isSegIntersect(Line<T> a, Line<U> b) {
	// 判断线段相交
	return CCW(b.p1, a) * CCW(b.p2, a) <= 0
	    && CCW(a.p1, b) * CCW(a.p2, b) <= 0;
}
		\end{minted}

		\textbf{命名提示}：函数名特意叫 \texttt{isSegIntersect}（不是 \texttt{isIntersect}）——\texttt{Line<T>} 在本模板里同时承担「直线」与「线段」两种语义，这里 \texttt{a, b} 是被当成线段判定的，函数名前加 \texttt{Seg} 防止与「两\textbf{直线}是否相交」（任意非平行直线都相交）混淆。

		\subsection{两条直线交点}

		\href{https://vjudge.net/solution/64186810}{vjudge 提交}。

		设两条直线分别用方向式 $p + t\vec{v}$、$q + s\vec{w}$。在交点处：
		\begin{align*}
			p + t\vec{v} &= q + s\vec{w} \\
			(p - q) \times \vec{w} + t\, \vec{v} \times \vec{w} &= 0 \quad (\text{两边叉乘 } \vec{w}\text{，自身叉乘消去 } s\vec{w}) \\
			t &= \dfrac{\vec{w} \times (q - p)}{\vec{v} \times \vec{w}}
		\end{align*}
		交点即 $p + t\vec{v}$。\textbf{返回值固定为 \texttt{Point<DB>}}（除法走浮点，不论输入 \texttt{T} 是什么都向 \texttt{DB} 收敛）。函数体内先把所有 \texttt{Point<T>} 走 promotion 提到 \texttt{Point<DB>}，整数 \texttt{T} 也合法（实参为整数 \texttt{Line<ll>} 时不再 \texttt{assert}）。

		\begin{minted}{cpp}
// 依赖：
//   struct Point, struct Line
//   cross(Point, Point)
// p+tv=q+sw
// (p-q) x w+tv x w = 0     (利用叉乘乘以自身为0，消掉sw)
// t=-(p-q)x w/(vxw)
// 传入的是方向式  拿到的是两直线交点
template<class T, class U>
Point<DB> getintersect(const Line<T> &a, const Line<U> &b) {
	Point<DB> p = a.p1, v = a.p2 - a.p1, q = b.p1, w = b.p2 - b.p1;
	Point<DB> u = p - q;
	DB t = (DB) cross(w, u) / cross(v, w);
	return p + v * t;
}
		\end{minted}

		\subsection{两线段间最小距离}

		\href{https://vjudge.net/solution/64187931}{vjudge 提交}。

		若两线段相交则距离为 $0$；否则取四个「端点 → 对方线段」距离的最小值。点到线段距离 \texttt{distancePS} 中要先判断投影是否落在线段两端之外，落到外面就退化为对应端点的距离。

		\begin{minted}{cpp}
// 依赖：
//   struct Point, struct Line
//   sgn, cross, dot
//   Point::abs()
//   isSegIntersect(Line, Line)
template<class T, class U>
DB distancePL(Point<T> p, Line<U> l) {
	auto val = cross(l.p1, l.p2, p); // 三参数：以 l.p1 为原点
	return sgn(val) * val / (l.p2 - l.p1).abs();
}

// AC https://qoj.ac/submission/2001668
template<class T, class U>
DB distancePS(Point<T> p, Line<U> l) {
	if (l.p1 == l.p2) return (p - l.p1).abs(); // 退化为点
	auto t = l.p2 - l.p1;
	if (sgn(dot(t, p - l.p1)) < 0) return (p - l.p1).abs();
	if (sgn(dot(t, p - l.p2)) > 0) return (p - l.p2).abs();
	return distancePL(p, l);
}

// AC https://qoj.ac/submission/2001588
template<class T, class U>
DB distanceSS(Line<T> a, Line<U> b) {
	if (isSegIntersect(a, b)) return 0;
	return min({distancePS(a.p1, b), distancePS(a.p2, b),
	            distancePS(b.p1, a), distancePS(b.p2, a)});
}
		\end{minted}

		\subsection{多边形面积}

		\href{https://vjudge.net/solution/64207564}{vjudge 提交}。

		用「向量叉乘求三角形有向面积之和」（鞋带公式）：
		$$
		A = \frac{1}{2} \left| \sum_{i=0}^{n-1} \mathrm{cross}(P_i, P_{(i+1) \bmod n}) \right|
		$$
		逆时针排列时 $\sum$ 本身已为正，顺时针为负，统一取绝对值。

		\textbf{拆出 \texttt{polygonSignedArea}}：返回值 \texttt{auto} 跟着 \texttt{cross} 走（整数 \texttt{T} $\to$ \texttt{i128}，浮点 \texttt{T} $\to$ \texttt{T}）。\textcolor{red}{\textbf{不要}}像旧版那样写成 \texttt{DB res = 0; res += cross(...)}——在整数 \texttt{T} 下 \texttt{cross} 返回 \texttt{\_\_int128\_t}，赋给 \texttt{double} 累加器会在 $|val| > 2^{53}$ 时丢精度。除此之外，\texttt{reorder\_polygon} 也复用 \texttt{polygonSignedArea} 取符号判方向，避免「前三点共线」时的局部退化。

		\begin{minted}{cpp}
// 依赖：
//   struct Point
//   cross(Point, Point)
//   Point::abs()

// 多边形 2 倍有向面积（不除 2 以保整数精度）
// 返回类型 auto：整数 T -> i128，浮点 T -> T
// > 0: CCW；< 0: CW；== 0: 退化（全共线 / 自交）
template<class T>
auto polygonSignedArea(const vector<Point<T>> &poly) {
	auto res = cross(Point<T>(), Point<T>()); // 零值，类型由 cross 继承
	size_t n = poly.size();
	for (size_t i = 0; i < n; ++i) res += cross(poly[i], poly[(i + 1) % n]);
	return res;
}

template<class T>
DB polygonArea(const vector<Point<T>> &poly) {
	auto s = polygonSignedArea(poly);
	if (s < 0) s = -s;
	return (DB) s / 2.0;
}

template<class T>
DB polygon_perimeter(const vector<Point<T>> &poly) {
	DB res = 0;
	size_t n = poly.size();
	for (int i = 0; i < n; ++i) res += (poly[i] - poly[(i + 1) % n]).abs();
	return res;
}

// 正 n 边形面积: A = n * l^2 / (4 * tan(pi / n)), n >= 3
DB regular_polygon_area(ll n, DB l) {
	return (DB) n * l * l / (4.0L * tanl(PI / (DB) n));
}
		\end{minted}

		\subsection{圆 \texttt{Circle}：测度五件套}

		圆心 \texttt{c} + 半径 \texttt{r} 的最小描述，配套\textbf{弧长 / 弦长 / 弓形高 / 扇形面积 / 弓形面积}五件套测度。

		\textbf{角度约定}：入参 \texttt{theta} 是弦所对的\textbf{圆心角}，单位\textbf{弧度}，约束 $\theta \in [0, 2\pi]$。$\theta < \pi$ 是 minor 弓形（小的那块），$\theta > \pi$ 是 major 弓形（大的那块），公式同一个，\textbf{不做内部 deg $\leftrightarrow$ rad 转换}（调用方自己 \texttt{deg * PI / 180}）。\texttt{point\_at} 接受任意 \texttt{theta}（含负值 / 超 $2\pi$），不抛 assert——绕圆走多圈合法。

		\textbf{弓形面积公式}：弓形 = 扇形 $-$ 三角形，化简得
		$$S_{\text{seg}}(\theta) = \tfrac{1}{2} r^2 (\theta - \sin \theta).$$
		AC 实战：\href{https://vjudge.net/problem/Gym-105161B}{Gym 105161B Area of the Devil}（圆 $-$ 五块 segment $-$ 五块三角形）。

		\textbf{\texttt{ccw\_angle} helper}：从极角 \texttt{theta\_from} CCW 转到 \texttt{theta\_to} 的弧度差，落在 $[0, 2\pi)$。\textcolor{red}{\textbf{不要}}用 \texttt{abs(theta\_to - theta\_from)}——后者在 wrap 处（例如 from $= 352^{\circ}$, to $= 54^{\circ}$）会算成 $298^{\circ}$ 而非真实的 $62^{\circ}$。

		\begin{minted}{cpp}
// 依赖：
//   struct Point, DB = long double, PI, eps
template<class T>
struct Circle {
	Point<T> c;
	T r;

	Circle() = default;
	Circle(Point<T> c_, T r_) : c(c_), r(r_) {}

	// Circle<U> -> Circle<T> 隐式 (common_type 单调升档方向, 同 Point converting ctor)
	template<class U, class = enable_if_t<is_same_v<common_type_t<U, T>, T>>>
	Circle(const Circle<U> &o) : c(o.c), r(o.r) {}

	// 圆上极角 theta 处的点 c + (r cos theta, r sin theta)
	// theta 不限范围（允许负值 / > 2 PI），始终走 cos/sin 不抛 assert
	Point<DB> point_at(DB theta) const {
		return Point<DB>((DB) c.x + (DB) r * cosl(theta),
		                 (DB) c.y + (DB) r * sinl(theta));
	}

	void check_theta(DB theta) const {
#ifdef LOCAL
		assert(theta >= -eps && theta <= 2 * PI + eps);
#endif
		(void) theta;
	}

	// 弧长 = r theta
	DB arc_length(DB theta) const {
		check_theta(theta);
		return (DB) r * theta;
	}

	// 弦长 = 2 r sin(theta / 2)
	DB chord_length(DB theta) const {
		check_theta(theta);
		return 2.0L * (DB) r * sinl(theta / 2);
	}

	// 弓形高 (sagitta) = r (1 - cos(theta / 2))
	DB sagitta(DB theta) const {
		check_theta(theta);
		return (DB) r * (1.0L - cosl(theta / 2));
	}

	// 扇形面积 = (1/2) r^2 theta
	DB sector_area(DB theta) const {
		check_theta(theta);
		return 0.5L * (DB) r * r * theta;
	}

	// 弓形面积 = (1/2) r^2 (theta - sin theta)
	// AC https://vjudge.net/solution/69814399 (Area of the Devil)
	DB segment_area(DB theta) const {
		check_theta(theta);
		return 0.5L * (DB) r * r * (theta - sinl(theta));
	}

	// AC https://vjudge.net/solution/69993234 使用 sgn 避免精度问题
	template<class U>
	bool is_point_in_circle(const Point<U> &p) {
		return sgn(distancePPNorm(p, c) - squared(r)) <= 0;
	}
};
		\end{minted}

		\subsection{两圆位置关系 \texttt{circle\_relation}}

		\href{https://vjudge.net/solution/69918813}{vjudge 提交}（\href{https://vjudge.net/problem/Aizu-CGL_7_A}{Aizu CGL\_7\_A}）。

		按圆心距 $d$ 与 $r_1 \pm r_2$ 的关系判 5 种位置：相离 / 外切 / 相交 / 内切 / 内含。\textbf{enum 取值即两圆公切线条数}（外离 4 / 外切 3 / 相交 2 / 内切 1 / 内含 0），方便直接拿来当答案输出。

		\textbf{重合圆 sentinel \texttt{COINCIDENT = -1}}：同心同径时公切线无穷多，单独一档；调用方按需特判（如题面要求输出 \texttt{-1} / \texttt{Inf} / 任意大数）。

		\textbf{用 $\mathrm{squared}(d)$ 与 $\mathrm{squared}(r_1 \pm r_2)$ 比较，不直接比 $d$}：整数情形 $(r_1 + r_2)^2$ 在 \texttt{ll} 边界（输入 1e9 量级）会溢，走 \texttt{squared()} 自动升 \texttt{\_\_int128\_t} 防溢出；浮点情形也省一次 \texttt{sqrt}。

		\begin{minted}{cpp}
// 依赖：
//   struct Circle, distancePPNorm, squared, sgn
enum CircleRelation {
	// 表示这个公切线个数; COINCIDENT (-1) 是重合 sentinel (无穷多公切线)
	COINCIDENT = -1,
	ONE_CONTAINS_OTHER = 0,
	INTERNALLY_TANGENT = 1,
	INTERSECTING = 2,
	EXTERNALLY_TANGENT = 3,
	EXTERNALLY_APART = 4,
};

template<class T>
CircleRelation circle_relation(const Circle<T> &a, const Circle<T> &b) {
	auto d2 = distancePPNorm(a.c, b.c);
	auto sum = squared(a.r + b.r);
	auto dif = squared(a.r - b.r);
	// 重合 (同心 + 同径): 无穷多公切线, 单独 sentinel
	if (sgn(d2) == 0 && sgn(a.r - b.r) == 0) return COINCIDENT;
	if (sgn(d2 - sum) > 0) return EXTERNALLY_APART;
	if (sgn(d2 - sum) == 0) return EXTERNALLY_TANGENT;
	if (sgn(d2 - dif) > 0) return INTERSECTING;
	if (sgn(d2 - dif) == 0) return INTERNALLY_TANGENT;
	return ONE_CONTAINS_OTHER;
}
		\end{minted}

		\subsection{三角形内接圆 / 外接圆}

		\textbf{内接圆 \texttt{incircle\_triangle}}：\href{https://vjudge.net/solution/69920478}{vjudge 提交}（\href{https://vjudge.net/problem/Aizu-CGL_7_B}{Aizu CGL\_7\_B}）。内心 $I$ 是三条\textbf{角平分线}的交点，用重心式直接写出：
		$$I = \frac{|BC| \cdot A + |AC| \cdot B + |AB| \cdot C}{|AB| + |AC| + |BC|}.$$
		内切圆半径 $r = $ 内心到任意一条边的距离，取 \texttt{min} 三条边距离做防御性数值兜底（三个值理论上相等，浮点下取最小避免 r 越界出三角形）。

		\textbf{外接圆 \texttt{outcircle\_triangle}}：\href{https://vjudge.net/solution/69927652}{vjudge 提交}（\href{https://vjudge.net/problem/Aizu-CGL_7_C}{Aizu CGL\_7\_C}）。外心 $O$ 是三条\textbf{垂直平分线}的交点，取 $AB$ 与 $AC$ 两边的中垂线求交即可：
		\begin{itemize}
			\setlength\itemsep{0pt}
			\item $AB$ 中点 $M_{AB} = (A + B) / 2$，中垂线方向 $= (B - A)$\texttt{.rot90()}；
			\item $AC$ 中点 $M_{AC} = (A + C) / 2$，中垂线方向 $= (C - A)$\texttt{.rot90()}；
			\item $O = \texttt{getintersect}(L_{AB}, L_{AC})$，$R = |O - A|$。
		\end{itemize}

		\textcolor{red}{\textbf{不要}}用「重心式」($I = (a \cdot BC + b \cdot AC + c \cdot AB)/(AB+AC+BC)$）求外心——那是\textbf{内心}公式，等腰 / 等边三角形样例下两者重合，会侥幸过样例但 WA。

		\textbf{退化 sentinel}：三点共线 / 任意两点重合（\texttt{cross(a, b, c) == 0}）时，内接圆 / 外接圆都不存在，分别返 \texttt{r = 0} 与 \texttt{r = -1} 作 sentinel；调用方用 \texttt{ci.r == 0} / \texttt{co.r < 0} 判退化即可，不会再吃到除零 NaN。

		\begin{minted}{cpp}
// 依赖：
//   struct Point, struct Line, struct Circle
//   distancePP, distancePS, getintersect, Point::rot90, sgn, cross（三参数）
// 退化 (共线 / 任意两点重合, cross == 0): 内接圆不存在, 返回 r=0 sentinel
Circle<DB> incircle_triangle(Point<DB> a, Point<DB> b, Point<DB> c) {
	if (sgn(cross(a, b, c)) == 0) return Circle<DB>(a, 0);
	DB ab = distancePP(a, b), ac = distancePP(a, c), bc = distancePP(b, c);
	Point<DB> I = (a * bc + ac * b + ab * c) / (ab + ac + bc);
	return Circle<DB>(I, min({distancePS(I, Line(a, c)),
	                          distancePS(I, Line(a, b)),
	                          distancePS(I, Line(b, c))}));
}

// 外心 = 两条中垂线交点
// 退化 (共线 / 任意两点重合, cross == 0): 外接圆不存在 (半径无穷大), 返回 r=-1 sentinel
Circle<DB> outcircle_triangle(Point<DB> a, Point<DB> b, Point<DB> c) {
	if (sgn(cross(a, b, c)) == 0) return Circle<DB>(a, -1);
	Point<DB> abMid = (a + b) / 2, acMid = (a + c) / 2;
	Point<DB> vecAb = (b - a).rot90(), vecAc = (c - a).rot90();
	Point<DB> I = getintersect(Line(abMid, abMid + vecAb),
	                           Line(acMid, acMid + vecAc));
	return Circle<DB>(I, distancePP(I, a));
}
		\end{minted}

		\subsection{直线与圆的关系 / 交点 \texttt{circle\_line\_relation} / \texttt{getCrossPointsCL}}

		\href{https://vjudge.net/solution/69928470}{vjudge 提交}（\href{https://vjudge.net/problem/Aizu-CGL_7_D}{Aizu CGL\_7\_D}）；弦长版 \href{https://vjudge.net/solution/69929454}{vjudge 提交}（\href{https://vjudge.net/problem/EOlymp-359}{EOlymp-359}）。

		\textbf{关系判定}：圆心 $C$ 到直线距离 $d = |\mathrm{cross}(p_1, p_2, C)| / |p_2 - p_1|$，与半径 $r$ 比较。两边平方 + 乘 $|p_2 - p_1|^2$ 走整数：$\mathrm{cross}^2$ vs $r^2 \cdot |p_2 - p_1|^2$，\texttt{squared()} 自动升 \texttt{\_\_int128\_t} 防溢。

		\textbf{交点求法}：先取垂足 $F = \mathrm{project}(C, l)$，半弦长 $h = \sqrt{r^2 - |CF|^2}$（斜边 $r$ 与直角边 $|CF|$ 在直角三角形里推另一直角边），单位方向 $\hat{d} = (p_2 - p_1) / |p_2 - p_1|$，交点 $F \pm \hat{d} \cdot h$。

		\textbf{退化 sentinel}：\texttt{l.p1 == l.p2} 时直接返 0 / 空集，不会跑到 \texttt{project} 的除零路径。

		\begin{minted}{cpp}
// 依赖：
//   struct Point, struct Line, struct Circle
//   cross（三参数）, squared, sgn, project, Point::abs
template<class T>
int circle_line_relation(const Circle<T> &cir, const Line<T> &l) {
	// 退化直线 (p1 == p2): 不构成有效直线, 视为无交点
	if (l.p1 == l.p2) return 0;
	auto cr2 = squared(cross(l.p1, l.p2, cir.c));
	auto lenSq = (l.p2 - l.p1).norm();
	auto rhs = squared(cir.r) * lenSq;
	int s = sgn(cr2 - rhs);
	if (s > 0) return 0;
	if (s == 0) return 1;
	return 2;
}

template<class T>
vector<Point<DB> > getCrossPointsCL(const Circle<T> &cir, const Line<T> &l) {
	int k = circle_line_relation(cir, l);
	if (k == 0) return {};
	// 得到垂足
	Point<DB> F = project(cir.c, l);
	if (k == 1) return {F};
	Point<DB> d = l.p2 - l.p1;
	Point<DB> dhat = d / d.abs();
	// 斜边 R 与直角边（圆心和垂足）得到另外一直角边（交点与垂足）
	DB h = sqrtl(max((DB)0, (DB)cir.r * cir.r - (cir.c - F).norm()));
	return {F - dhat * h, F + dhat * h};
}
		\end{minted}

		\subsection{两圆的关系 / 交点 \texttt{circle\_circle\_relation} / \texttt{getCrossPointsCC}}

		\href{https://vjudge.net/solution/69929473}{vjudge 提交}（\href{https://vjudge.net/problem/Aizu-CGL_7_E}{Aizu CGL\_7\_E}）；含重合 sentinel 版 \href{https://vjudge.net/solution/69929406}{vjudge 提交}（\href{https://vjudge.net/problem/EOlymp-4779}{EOlymp-4779}）。

		\textbf{关系判定}：$d^2$、$(r_1 + r_2)^2$、$(r_1 - r_2)^2$ 全走整数 \texttt{squared} 比较，不开方。$d^2 > (r_1 + r_2)^2$ 相离；$d^2 = (r_1 + r_2)^2$ 外切；$(r_1 - r_2)^2 < d^2 < (r_1 + r_2)^2$ 相交；$d^2 = (r_1 - r_2)^2$ 内切；$d^2 < (r_1 - r_2)^2$ 内含。

		\textbf{交点求法}：连心线方向 $\hat{d} = (c_2.c - c_1.c) / D$。弦中点 $M = c_1.c + \hat{d} \cdot a$，其中 $a = (r_1^2 - r_2^2 + D^2) / (2 D)$（联立两圆方程相减消平方项得到）。沿连心线垂直方向 $\hat{d}.\mathrm{rot90}()$ $\pm$ 半弦长 $h = \sqrt{r_1^2 - a^2}$ 即两交点。

		\textbf{重合圆 sentinel}：同心同径时 \texttt{circle\_circle\_relation} 返 $-1$（交点无穷多），\texttt{getCrossPointsCC} 据此返空，重合语义由 caller 自己处理。

		\textbf{promotion 坑}：\texttt{(c2.c - c1.c) / D} 若 \texttt{cN} 是 \texttt{Point<ll>}、$D$ 是 \texttt{DB}，会走 \texttt{Point<ll>::operator/(Point<ll>, ll)} 把 $D$ 截成整数。必须先 \texttt{Point<DB> d = c2.c - c1.c;}。

		\begin{minted}{cpp}
// 依赖：
//   struct Point, struct Circle
//   distancePPNorm, squared, sgn, Point::abs, Point::rot90
template<class T>
int circle_circle_relation(const Circle<T> &a, const Circle<T> &b) {
	auto d2 = distancePPNorm(a.c, b.c);
	auto sum = squared(a.r + b.r);
	auto dif = squared(a.r - b.r);
	// 重合 (同心 + 同径): 交点无穷多, 返 -1 sentinel
	if (sgn(d2) == 0 && sgn(a.r - b.r) == 0) return -1;
	if (sgn(d2 - sum) > 0) return 0;
	if (sgn(d2 - dif) < 0) return 0;
	if (sgn(d2 - sum) == 0 || sgn(d2 - dif) == 0) return 1;
	return 2;
}

template<class T>
vector<Point<DB> > getCrossPointsCC(const Circle<T> &c1, const Circle<T> &c2) {
	int k = circle_circle_relation(c1, c2);
	// k = -1 重合 (无穷多) / k = 0 不交: 统一返空, 重合语义由 caller 自己处理
	if (k <= 0) return {};
	// 弦中点
	Point<DB> d = c2.c - c1.c;
	DB D = d.abs();
	Point<DB> dhat = d / D;
	DB a = ((DB)c1.r * c1.r - (DB)c2.r * c2.r + D * D) / (2 * D);
	Point<DB> M = (Point<DB>)c1.c + dhat * a;
	if (k == 1) return {M};
	// 沿连心线垂直方向 +- 半弦长
	DB h = sqrtl(max((DB)0, (DB)c1.r * c1.r - a * a));
	Point<DB> nhat = dhat.rot90();
	return {M + nhat * h, M - nhat * h};
}
		\end{minted}

		\subsection{过一点圆的切点 \texttt{getTangentPoints}}

		\href{https://vjudge.net/solution/69929512}{vjudge 提交}（\href{https://vjudge.net/problem/Aizu-CGL_7_F}{Aizu CGL\_7\_F}）。

		\textbf{推导}：过外部点 $P$ 作圆 $(C, r)$ 切线时，切点 $T$ 同时满足 $PT \perp CT$ 与 $|CT| = r$，即 $\triangle PCT$ 在 $T$ 处是直角三角形，$|PT| = \sqrt{|PC|^2 - r^2}$。切点 $T$ 同时落在原圆 $(C, r)$ 与辅圆 $(P, |PT|)$ 上，对两圆求交即得。

		\textbf{三种退化分类直接 early return}：$P$ 在圆内（含圆心）返空；$P$ 在圆上唯一切点就是 $P$ 自身；$P$ 在圆外才走辅圆链。比"全部喂给 \texttt{getCrossPointsCC} 让它兜底"显式可读，也免依赖辅圆 r = 0 时的边界稳定性。

		\begin{minted}{cpp}
// 依赖：
//   struct Point, struct Circle
//   getCrossPointsCC
// 退化: P 在圆内 (含圆心) 返空; P 在圆上唯一切点就是 P 自身
vector<Point<DB> > getTangentPoints(Point<ll> P, Circle<ll> cir) {
	auto d2 = (P - cir.c).norm();
	auto r2 = (__int128_t)cir.r * cir.r;
	if (d2 < r2) return {};
	if (d2 == r2) return { Point<DB>((DB)P.x, (DB)P.y) };
	DB Lsq = (DB)d2 - (DB)r2;
	DB L = sqrtl(max((DB)0, Lsq));
	Circle<DB> orig((Point<DB>)cir.c, (DB)cir.r);
	Circle<DB> aux((Point<DB>)P, L);
	return getCrossPointsCC(orig, aux);
}
		\end{minted}

		\subsection{两圆公切线 \texttt{getCommonTangentPoints}}

		\href{https://vjudge.net/solution/69929554}{vjudge 提交}（\href{https://vjudge.net/problem/Aizu-CGL_7_G}{Aizu CGL\_7\_G}）。函数返回的是公切线在 $c_1$ 上的切点列表（每条公切线一个点；对应 $c_2$ 上的切点由切线本身确定）。

		\textbf{推导}：设 $T_1 = c_1.c + r_1 \cdot \vec{n}$，$\vec{n}$ 是单位法向量。令 $\vec{u} = (c_2.c - c_1.c) / D$，分外 / 内公切两类：
		\begin{itemize}
			\setlength\itemsep{0pt}
			\item 外公切：两边法向同向，$T_2 = c_2.c + r_2 \cdot \vec{n}$ 且 $(T_2 - T_1) \perp \vec{n}$，得 $(c_2.c - c_1.c) \cdot \vec{n} = r_1 - r_2$，故 $\cos \alpha = (r_1 - r_2) / D$；
			\item 内公切：两边法向反向，$T_2 = c_2.c - r_2 \cdot \vec{n}$，得 $\cos \alpha = (r_1 + r_2) / D$。
		\end{itemize}
		然后 $\vec{n} = \cos \alpha \cdot \vec{u} + \sin \alpha \cdot \vec{u}_\perp$，$\sin \alpha = \pm \sqrt{1 - \cos^2 \alpha}$ 取两组。

		\textbf{退化}：$|\cos \alpha| > 1$ 该类无切（内含 / 包含跳过）；$\sin \alpha = 0$ 该类只 1 条（内 / 外切临界）；$D = 0$（同心，含重合）统一返空。最后 \texttt{sort + unique} 去重，应对内 / 外切临界 + $r = 0$ 退化下 4 个候选近重合的情形。

		\begin{minted}{cpp}
// 依赖：
//   struct Point, struct Circle, eps
//   Point::abs, Point::rot90
vector<Point<DB> > getCommonTangentPoints(Circle<ll> c1, Circle<ll> c2) {
	vector<Point<DB> > ret;
	Point<DB> d = c2.c - c1.c;
	DB D = d.abs();
	if (D < eps) return ret;
	Point<DB> u = d / D;
	Point<DB> u_perp = u.rot90();
	DB r1 = c1.r, r2 = c2.r;
	auto add = [&](DB cos_a) {
		if (cos_a > 1 + eps || cos_a < -1 - eps) return;
		DB s = sqrtl(max((DB)0, 1 - cos_a * cos_a));
		Point<DB> c1d = c1.c;
		if (s < eps) {
			Point<DB> n = u * cos_a;
			ret.push_back(c1d + n * r1);
		} else {
			ret.push_back(c1d + (u * cos_a + u_perp * s) * r1);
			ret.push_back(c1d + (u * cos_a - u_perp * s) * r1);
		}
	};
	add((r1 - r2) / D); // 外公切
	add((r1 + r2) / D); // 内公切
	sort(ret.begin(), ret.end());
	// 内/外切临界 + r=0 退化时 4 候选可能近重合, unique 去重
	ret.erase(unique(ret.begin(), ret.end()), ret.end());
	return ret;
}
		\end{minted}

		\subsection{两圆交集面积 \texttt{two\_circle\_intersect\_area}}

		\href{https://vjudge.net/solution/69929585}{vjudge 提交}（\href{https://vjudge.net/problem/Aizu-CGL_7_I}{Aizu CGL\_7\_I}）。

		\textbf{推导}：两圆相交时公共弦把每个圆切成两个弓形（segment / cap）。交集面积 = $c_1$ 靠 $c_2$ 那侧的弓形 + $c_2$ 靠 $c_1$ 那侧的弓形。圆心角由余弦定理给出：
		$$\theta_1 = 2 \arccos \frac{D^2 + r_1^2 - r_2^2}{2 D r_1}, \quad \theta_2 = 2 \arccos \frac{D^2 + r_2^2 - r_1^2}{2 D r_2}.$$
		单个弓形面积 $= \tfrac{1}{2} r^2 (\theta - \sin \theta)$。

		\textbf{边界}：$D \ge r_1 + r_2$ 不交 $\to 0$；$D \le |r_1 - r_2|$ 一圆完全包含另一圆 $\to \pi \cdot r_{\min}^2$（重合圆自然落入此档，无需特判）。

		\begin{minted}{cpp}
// 依赖：
//   struct Circle
//   distancePP, sgn, PI
template<class T>
DB two_circle_intersect_area(Circle<T> c1, Circle<T> c2) {
	DB D = distancePP(c1.c, c2.c);
	DB r1 = c1.r, r2 = c2.r;
	if (sgn(D - r1 - r2) >= 0) return 0;
	DB diff = abs(r1 - r2);
	if (sgn(D - diff) <= 0) {
		DB rmin = min(r1, r2);
		return PI * rmin * rmin;
	}
	DB t1 = 2 * acosl((D * D + r1 * r1 - r2 * r2) / (2 * D * r1));
	DB t2 = 2 * acosl((D * D + r2 * r2 - r1 * r1) / (2 * D * r2));
	return 0.5L * r1 * r1 * (t1 - sinl(t1)) + 0.5L * r2 * r2 * (t2 - sinl(t2));
}
		\end{minted}

		\subsection{简单多边形与圆的交集面积 \texttt{polygon\_circle\_intersect\_area}}

		\href{https://vjudge.net/solution/69929628}{vjudge 提交}（\href{https://vjudge.net/problem/Aizu-CGL_7_H}{Aizu CGL\_7\_H}）。

		\textbf{核心}：把多边形按边累加「三角形 $(O, A_i, A_{i+1})$ 与 disk $(O, r)$ 的有向面积」，$O$ 是圆心。每条边的贡献由 \texttt{tri\_disk\_signed\_area} 计算：先把 $A, B$ 平移到以 $O$ 为原点的坐标系；然后按 $A, B$ 是否在圆内分四种 case：
		\begin{itemize}
			\setlength\itemsep{0pt}
			\item 两端都在圆内：整段就是真三角形，有向面积 $= \tfrac{1}{2} \mathrm{cross}(A, B)$；
			\item 两端都在圆外但弦穿圆得入 / 出两点 $P_1, P_2$：扇形 $OAP_1$ + 真三角形 $OP_1 P_2$ + 扇形 $OP_2 B$；
			\item 两端都在圆外且不穿圆：整段对应一个扇形（从 $\vec{OA}$ 转到 $\vec{OB}$）；
			\item $A$ 内 $B$ 外（或反之）：在边上找穿圆点 $P$，$OAP$ 是真三角形 + $OPB$ 是扇形。
		\end{itemize}
		最后整圈累加取 \texttt{abs}，输入多边形 CW / CCW 都可。\textbf{扇形角度用 \texttt{atan2(cross, dot)} 算}：能拿到有向角度 $\in (-\pi, \pi]$，正负自动包含弧的方向（不像 \texttt{acos} 只给 $[0, \pi]$ 丢方向信息）。

		\textbf{退化}：\texttt{tri\_disk\_signed\_area} 在 $A = B$（多边形相邻顶点重合）时直接返 0，整体也吃得到。

		\begin{minted}{cpp}
// 依赖：
//   struct Point, struct Circle
//   cross（两参数）, dot, Point::abs, sgn, sqrtl, atan2l
DB tri_disk_signed_area(Point<DB> A, Point<DB> B, DB r) {
	if (A == B) return 0; // 零长度边, 三角形退化, 有向面积为 0
	DB rA = A.abs(), rB = B.abs();
	bool aIn = sgn(rA - r) <= 0;
	bool bIn = sgn(rB - r) <= 0;
	auto ang = [](Point<DB> P, Point<DB> Q) {
		return atan2l((DB)cross(P, Q), (DB)dot(P, Q));
	};
	if (aIn && bIn) {
		// 整段都在圆内, 真三角形面积
		return 0.5L * (DB)cross(A, B);
	}
	// 求 A + t*(B-A) 落圆上的 t: |A + t d|^2 = r^2
	Point<DB> d = B - A;
	DB dd = (DB)dot(d, d);
	DB dA = (DB)dot(d, A);
	DB AA = (DB)dot(A, A);
	DB disc = dA * dA - dd * (AA - r * r);
	disc = max(disc, (DB)0);
	DB sq = sqrtl(disc);
	DB t1 = (-dA - sq) / dd, t2 = (-dA + sq) / dd;
	if (!aIn && !bIn) {
		// 两端都在外, 看弦区间 [t1, t2] 与段 [0, 1] 是否相交
		if (sgn(t2) < 0 || sgn(t1 - 1) > 0 || sgn(disc) <= 0) {
			// 不穿圆, 整段对应一个扇形 (从 OA 转到 OB)
			return 0.5L * r * r * ang(A, B);
		}
		t1 = max(t1, (DB)0);
		t2 = min(t2, (DB)1);
		Point<DB> P1 = A + d * t1;
		Point<DB> P2 = A + d * t2;
		return 0.5L * r * r * ang(A, P1) + 0.5L * (DB)cross(P1, P2)
		     + 0.5L * r * r * ang(P2, B);
	}
	if (aIn && !bIn) {
		// A 内 B 外, 段 [0,1] 在 [t1, t2] 内的部分: t >= 0 (因为 A 内 t1 < 0), 取 t2
		Point<DB> P = A + d * t2;
		return 0.5L * (DB)cross(A, P) + 0.5L * r * r * ang(P, B);
	}
	// !aIn && bIn: A 外 B 内, t1 是入口
	Point<DB> P = A + d * t1;
	return 0.5L * r * r * ang(A, P) + 0.5L * (DB)cross(P, B);
}

// 简单多边形 (CCW) 与圆的交集面积. 累加每条边的 tri_disk_signed_area.
DB polygon_circle_intersect_area(Circle<DB> cir, const vector<Point<ll> > &poly) {
	DB res = 0;
	int n = poly.size();
	for (int i = 0; i < n; ++i) {
		Point<DB> A = poly[i] - cir.c;
		Point<DB> B = poly[(i + 1) % n] - cir.c;
		res += tri_disk_signed_area(A, B, cir.r);
	}
	return abs(res);
}
		\end{minted}

		\subsection{多边形凸性检验}

		\textbf{允许的输入}：\textcolor{red}{\textbf{不允许自交}}的多边形，顺序可顺可逆，函数会顺手把 \texttt{poly} 旋转 / 反向到「最左下点位 \texttt{[0]} + CCW 顺序」。

		下面是三类\textbf{不被认为是凸多边形}的反例。

		\subsubsection{反例 1：内角为反向凹角（外角为锐角）}

		外角是锐角时 $\sin > 0$，叉乘也 $> 0$，看似 CCW；但其实形成的是\textcolor{red}{凹陷}。

		\begin{center}
			\begin{tikzpicture}[>=Stealth]
				\draw[thick] (0,0) -- (3,0) -- (3,2) -- (1.5,1) -- (0,2) -- cycle;
				\fill[red] (1.5,1) circle (2pt);
				\node[font=\scriptsize, color=red, anchor=west] at (1.6,1) {反向凹角};
			\end{tikzpicture}

			\vspace{2pt}
			{\footnotesize\sffamily\color{gray} 图注：标红顶点的内角 $> 180^{\circ}$（反射角），\texttt{isConvex} 中第一遍 \texttt{cross(poly[i] - poly[0], poly[i+1] - poly[0]) <= 0} 触发返回 false。}
		\end{center}

		\subsubsection{反例 2：自交多边形（蝴蝶结）}

		\begin{center}
			\begin{tikzpicture}[>=Stealth]
				\draw[thick] (0,0) -- (3,2) -- (3,0) -- (0,2) -- cycle;
				\fill[red] (1.5,1) circle (2pt);
				\node[font=\scriptsize, color=red, anchor=west] at (1.6,1) {自交点};
			\end{tikzpicture}

			\vspace{2pt}
			{\footnotesize\sffamily\color{gray} 图注：四个顶点排成蝴蝶结，相对 \texttt{poly[0]} 不再是单调 CCW 排列，第一遍叉乘扫描即被排除。}
		\end{center}

		\subsubsection{反例 3：内角恰为 $180^{\circ}$}

		三点共线（叉乘 \texttt{= 0}）也不允许，否则会让边数虚假增加。

		\begin{center}
			\begin{tikzpicture}[>=Stealth]
				\draw[thick] (0,0) -- (1.5,0) -- (3,0) -- (1.5,2.2) -- cycle;
				\fill[red] (1.5,0) circle (2pt);
				\node[font=\scriptsize, color=red, anchor=north] at (1.5,-0.05) {$180^{\circ}$};
			\end{tikzpicture}

			\vspace{2pt}
			{\footnotesize\sffamily\color{gray} 图注：底边上多了一个 $180^{\circ}$ 顶点，第二遍 \texttt{cross(a, b) <= 0} 等号触发返回 false（严格大于零才算凸）。}
		\end{center}

		\textbf{实现}：\texttt{reorder\_polygon} 把多边形旋转 / 反向到「最左下点位 \texttt{[0]} + CCW 顺序」，再分两轮检查：(a)~所有顶点相对于 \texttt{poly[0]} 是 CCW 排列；(b)~每个内角的叉乘符号合规。

		\textbf{\texttt{reorder\_polygon} 用整体 \texttt{polygonSignedArea} 判方向}，\textcolor{red}{\textbf{不要}}用「前三点 \texttt{cross}」局部判向：含三点共线的多边形（如长方形 + 边密集共线）前三点 \texttt{cross = 0}，会被旧版 \texttt{<= 0} 误判为 CW 而强 reverse，下游 \texttt{isConvex} / 旋转卡壳全部失效（QOJ 784 实测）。

		\textbf{\texttt{isConvex} 加 \texttt{strict} 参数}：
		\begin{itemize}
					\setlength\itemsep{0pt}
				\item \texttt{strict = true}（默认）：严格凸，拒绝三点共线 / 内角 = $180^{\circ}$；
				\item \texttt{strict = false}：弱凸，允许内角 = $180^{\circ}$（题面用「凸」时常允许共线，要的就是这个口径）。
		\end{itemize}
		\textbf{加宽接受范围有时是发现底层 bug 的好工具}：本次就是因为加 \texttt{strict = false} 后 lax 模式 fail，才暴露出 \texttt{reorder\_polygon} 的隐藏 bug——strict 模式因为 \texttt{<= 0} 提前 short-circuit 把它掩盖了。

		\begin{minted}{cpp}
// 依赖：
//   struct Point
//   sgn, cross（含三参数）, polygonSignedArea
template<class T>
void reorder_polygon(vector<Point<T>> &poly) {
	ll n = poly.size();
	// 整体有向面积判向，避免「前三点共线」的局部退化（QOJ 784 反例）
	if (sgn(polygonSignedArea(poly)) < 0) reverse(poly.begin(), poly.end());
	// 找到 y 最小、再 x 最小的点并循环左移到 [0]
	ll pos_mn = 0;
	for (int i = 1; i < n; ++i) if (poly[i].y < poly[pos_mn].y || (poly[i].y == poly[pos_mn].y && poly[i].x < poly[pos_mn].x)) pos_mn = i;
	rotate(poly.begin(), poly.begin() + pos_mn, poly.end());
}

// 不允许自交的多边形，顺序都可以；会把 poly 变为 CCW 顺序
// AC https://qoj.ac/submission/2002121
// strict = true : 严格凸（拒绝三点共线 / 内角 = 180°）
// strict = false: 弱凸  （允许三点共线 / 内角 = 180°），题面允许共线时用
template<class T>
bool isConvex(vector<Point<T>> poly, bool strict = true) {
	ll n = poly.size();
	if (n <= 2) return false; // 0 / 1 / 2 点不是多边形
	reorder_polygon(poly);
	auto fail = [&](int s) { return strict ? s <= 0 : s < 0; };
	// 所有顶点相对 poly[0] 应按 CCW 整齐排列，否则是星形 / 乱序
	bool has_turn = false; // 至少一处真左转，排除全共线（面积 = 0）退化
	for (int i = 1; i < n - 1; ++i) {
		int s = sgn(cross(poly[0], poly[i], poly[i + 1]));
		if (fail(s)) return false;
		if (s > 0) has_turn = true;
	}
	// 弱凸下若全程 cross == 0（全部点共线、面积 = 0）也得拒；
	// strict 下 has_turn 在第一轮必为 true，此判等价空操作
	if (!has_turn) return false;
	// strict 时拒绝内角 = 180°（sgn = 0），弱凸只拒绝凹（sgn < 0）
	for (int i = 0; i < n; ++i) {
		ll i1 = (i + 1) % n, i2 = (i + 2) % n;
		auto a = poly[i1] - poly[i], b = poly[i2] - poly[i1];
		if (fail(sgn(cross(a, b)))) return false;
	}
	return true;
}
		\end{minted}

		\subsection{点与多边形关系}

		多边形未必凸——「有向面积一定为正」的方法在凹多边形上失效，所以采用经典的\textbf{射线奇偶法}。

		\begin{center}
			\begin{tikzpicture}[>=Stealth]
				\draw[thick] (0,0) -- (3.2,0.3) -- (3.6,2) -- (2.2,1.4) -- (2,2.7) -- (0.5,2.2) -- cycle;
				\coordinate (Pin)  at (1.4, 1.4);
				\coordinate (Pout) at (1.4, -0.6);
				\fill[red]   (Pin)  circle (2pt) node[above, font=\scriptsize, color=red]   {$P_{\text{in}}$};
				\fill[blue]  (Pout) circle (2pt) node[below, font=\scriptsize, color=blue]  {$P_{\text{out}}$};
				\draw[->, dashed, red]  (Pin)  -- (5.0, 1.4);
				\draw[->, dashed, blue] (Pout) -- (5.0, -0.6);
				\node[font=\tiny, color=red,  anchor=south] at (4.0, 1.45) {1 次相交（奇）→ 在内};
				\node[font=\tiny, color=blue, anchor=north] at (4.0, -0.65) {0 次相交（偶）→ 在外};
			\end{tikzpicture}

			\vspace{2pt}
			{\footnotesize\sffamily\color{gray} 图注：从待判定点向任意方向打一条射线，统计与多边形边的相交次数。\textbf{奇 = 在内，偶 = 在外}。代码里射线方向用 \texttt{rng()} 随机，避开「射线恰好穿过顶点」的共线退化。}
		\end{center}

		返回值约定：\texttt{2}~=~严格在内、\texttt{1}~=~点落在某条边上、\texttt{0}~=~严格在外。

		\begin{minted}{cpp}
// 依赖：
//   struct Point, struct Line
//   CCW(Point, Line), ON_SEGMENT
//   isSegIntersect(Line, Line)
//   mt19937_64 rng（基础模板已定义）
//   宏 SZ（用户自定义，等价于 (int)x.size()）
template<class T>
int isInPolygon(Point<T> p, const vector<Point<T>> &poly) {
	ll res = 0;
	Line<T> l(p, p + Point<T>(2e18, rng()));
	for (int i = 0; i < poly.size(); ++i) {
		int i1 = (i + 1) % SZ(poly);
		Line<T> seg = Line<T>(poly[i], poly[i1]);
		int ccw = CCW(p, seg);
		if (ccw == ON_SEGMENT) return 1;
		if (isSegIntersect(l, seg)) ++res;
	}
	// 奇偶法则：射线相交奇数条边 = 在内
	return (res & 1) ? 2 : 0;
}
		\end{minted}

		\subsection{极角排序}

		按 \texttt{atan2(y, x)} 排序时，注意 \texttt{atan2} 浮点本身的精度问题。\textbf{规避调用 atan2} 的做法：先按「下半平面 / x 轴正向 / 上半平面」三档分桶，桶内再用叉乘符号判定相对方位。这样所有比较都用整数叉乘完成，没有浮点误差。

		起点是 $x$ 轴正半轴（含原点），逆时针环绕一周到 $x$ 轴负半轴（不含）。

		\begin{minted}{cpp}
// 依赖：
//   struct Point
//   sgn, cross
template<class T>
int get_part_of_point(const Point<T> &p) {
	if (p.y < 0) return 0;             // 下半平面 (-pi, 0)
	if (p.y == 0 && p.x >= 0) return 1; // 角度为 0
	return 2;                           // 上半平面 (0, pi]
}

// 按 atan2(y, x) 升序排，但用整数叉乘规避浮点
// AC https://qoj.ac/submission/2001485
template<class T>
void polar_angle_sort(vector<Point<T>> &poly) {
	sort(poly.begin(), poly.end(), [](const Point<T> &a, const Point<T> &b) -> bool {
		int part1 = get_part_of_point(a), part2 = get_part_of_point(b);
		if (part1 != part2) return part1 < part2;
		return sgn(cross(a, b)) == 1;
	});
}

// 重载：对有向直线按方向向量做极角排序（半平面交用）
// 同向（平行）的两条线，取「半平面内侧更靠右」的那条排前面——
// 等价判定：bl.p1 在 al 右侧（CLOCKWISE）时 al 在前
template<class T>
void polar_angle_sort(vector<Line<T>> &poly) {
	sort(poly.begin(), poly.end(), [](const Line<T> &al, const Line<T> &bl) -> bool {
		auto a = al.p2 - al.p1, b = bl.p2 - bl.p1;
		int part1 = get_part_of_point(a), part2 = get_part_of_point(b);
		if (part1 != part2) return part1 < part2;
		int c = sgn(cross(a, b));
		if (c != 0) return c == 1;
		// 同向：bl 在 al 右侧，说明 al 的半平面把 bl 完全包住，al 更紧、应排前
		return CCW(bl.p1, al) == CLOCKWISE;
	});
}

// polar_angle_sort 之后，把同向平行线去重，只保留每组最紧的那一条
// （配合上面的排序顺序，每组同向线的第一条即最紧）
template<class T>
void unique_line_polar_angle_sort(vector<Line<T>> &poly) {
	polar_angle_sort(poly);
	poly.resize(unique(all(poly), [](const Line<T> &al, const Line<T> &bl) {
		auto a = al.p2 - al.p1, b = bl.p2 - bl.p1;
		return sgn(cross(a, b)) == 0;
	}) - poly.begin());
}
		\end{minted}

		\subsection{Minkowski 闵可夫斯基和}

		给定两个凸包 $P_1$、$P_2$，定义其闵可夫斯基和：
		$$
		P_1 \oplus P_2 = \{a + b \mid a \in P_1,\ b \in P_2\}
		$$
		结果仍是凸包。\textbf{减法}就等价于「加上对方关于原点的中心对称图形」，这是判断两凸包是否相交、求最近距离的常用技巧。

		\begin{center}
			\begin{tikzpicture}[>=Stealth]
				\begin{scope}
					\draw[blue, thick] (-0.6,-0.6) rectangle (0.6, 0.6);
					\node[font=\scriptsize] at (0, -1.0) {$P_1$};
				\end{scope}
				\begin{scope}[xshift=2.0cm]
					\node[font=\scriptsize, anchor=center] at (-0.85, 0) {$\oplus$};
					\draw[orange, thick] (0,0) circle (0.4);
					\node[font=\scriptsize] at (0, -1.0) {$P_2$};
				\end{scope}
				\begin{scope}[xshift=4.4cm]
					\node[font=\scriptsize, anchor=center] at (-0.95, 0) {$=$};
					\draw[teal, thick] (-0.6, -1.0) -- (0.6, -1.0);
					\draw[teal, thick] (-0.6,  1.0) -- (0.6,  1.0);
					\draw[teal, thick] (-1.0, -0.6) -- (-1.0, 0.6);
					\draw[teal, thick] ( 1.0, -0.6) -- (1.0,  0.6);
					\draw[teal, thick] (0.6, 1.0)  arc (90:0:0.4);
					\draw[teal, thick] (1.0, -0.6) arc (0:-90:0.4);
					\draw[teal, thick] (-0.6, -1.0) arc (270:180:0.4);
					\draw[teal, thick] (-1.0, 0.6)  arc (180:90:0.4);
					\node[font=\scriptsize] at (0, -1.4) {$P_1 \oplus P_2$};
				\end{scope}
			\end{tikzpicture}

			\vspace{2pt}
			{\footnotesize\sffamily\color{gray} 图注：正方形 + 圆 = 圆角矩形——四条直边长度等于原正方形的边长，四个角是半径等于圆半径的四分之一圆弧（这是几何上严格的 Minkowski 和）。}
		\end{center}

		\textbf{算法核心}：先把两个凸包都 \texttt{reorder\_polygon}（最左下点起、CCW 顺序），然后\textbf{把它们的边向量按极角合并}（类似归并排序的归并步）：每次比较两条候选边向量的叉乘符号，挑极角更小的那条接到答案末尾。起点设为 $P_1[0] + P_2[0]$，沿合并后的边向量逐条累加即得新凸包。

		\begin{minted}{cpp}
// 依赖：
//   struct Point
//   sgn, cross
//   reorder_polygon（见多边形凸性检验一节）
// 参考 hh2048/XCPC 的 mincowski    AC https://qoj.ac/submission/2002545
template<class T>
vector<Point<T>> minkowski(vector<Point<T>> &P1, vector<Point<T>> &P2) {
	reorder_polygon(P1);
	reorder_polygon(P2);
	int n = P1.size(), m = P2.size();
	vector<Point<T>> V1(n), V2(m);
	for (int i = 0; i < n; i++) V1[i] = P1[(i + 1) % n] - P1[i];
	for (int i = 0; i < m; i++) V2[i] = P2[(i + 1) % m] - P2[i];
	vector<Point<T>> ans = {P1.front() + P2.front()};
	int i = 0, j = 0;
	while (i < n && j < m) {
		// cross(V1, V2) > 0 说明 V1 极角更小（CCW 下），先选小的
		Point<T> val = sgn(cross(V1[i], V2[j])) > 0 ? V1[i++] : V2[j++];
		ans.push_back(ans.back() + val);
	}
	while (i < n) ans.push_back(ans.back() + V1[i++]);
	while (j < m) ans.push_back(ans.back() + V2[j++]);
	return ans;
}
		\end{minted}

		\subsection{半平面交（half\_plane\_cut）}

		把每条有向直线 \texttt{Line(p1, p2)} 看作半平面 ——「沿 \texttt{p1 → p2} 方向走，左手边」是合法区域。所有半平面的交集若非空则为一个凸多边形（可能退化为无界、零面积等情形，本模板返回 \texttt{vector<Point<DB>>}，空表示交集为空 / 退化）。

		\textbf{算法核心}：双端队列扫描法。先 \texttt{unique\_line\_polar\_angle\_sort} 把直线按极角排序 + 同向去重，然后逐条加入 deque \texttt{q}：每次新直线 \texttt{l} 进来，先用 \texttt{l} 反向清理队尾（队尾两条线的交点若在 \texttt{l} 右侧，说明该交点被 \texttt{l} 切掉，对应的队尾线已经无贡献），再清理队首（对称地）；扫完一轮后还要做一次 \textbf{wrap-around}：用队首 / 队尾互相再清一遍（处理环形闭合时残留的失效线段）。

		\textbf{复杂度}：$O(N \log N)$（瓶颈是极角排序，扫描本身均摊 $O(N)$）。

		\textbf{常见用法}：
		\begin{itemize}
			\item \textbf{多个半平面求公共凸区域}：每个不等式 $a_i x + b_i y \le c_i$ 表示一个半平面，加入对应的有向直线即可。
			\item \textbf{多边形并的总面积（旋转切割型）}：经典题 \href{https://vjudge.net/problem/OpenJ_Bailian-2451}{\emph{Uyuw's Concert}}（多次直线切凸多边形求剩余面积）/ \href{https://www.luogu.com.cn/problem/P4196}{洛谷 P4196 [CQOI2006] 凸多边形}（多个凸多边形求交）。
			\item \textbf{有界化}：当输入只有半平面、没有边界 box 时，可手动加 4 条「大正方形边界」直线（如 \texttt{[-1e9, 1e9]} 见方），避免交集无界。
		\end{itemize}

		\begin{minted}{cpp}
// 依赖：
//   struct Point, struct Line
//   CCW, CLOCKWISE
//   getintersect, polygonArea
//   unique_line_polar_angle_sort
// AC https://vjudge.net/solution/69843663（Uyuw's Concert）
template<class T>
vector<Point<DB>> half_plane_cut(vector<Line<T>> lines) {
	unique_line_polar_angle_sort(lines);
	deque<Line<T>> q;
	for (auto l : lines) {
		if (SZ(q) < 2) { q.push_back(l); continue; }
		// 清队尾：队尾两条线交点若在 l 右侧（被 l 切掉），队尾线作废
		auto checkBack = [&]() {
			auto p = getintersect(q.rbegin()[0], q.rbegin()[1]);
			return CCW(p, l) != CLOCKWISE;
		};
		while (SZ(q) >= 2 && !checkBack()) q.pop_back();
		// 对称清队首
		auto checkFront = [&]() {
			auto p = getintersect(q.begin()[0], q.begin()[1]);
			return CCW(p, l) != CLOCKWISE;
		};
		while (SZ(q) >= 2 && !checkFront()) q.pop_front();
		q.push_back(l);
	}
	// wrap-around：用队首线再清一遍队尾
	while (SZ(q) >= 3) {
		auto p = getintersect(q.rbegin()[0], q.rbegin()[1]);
		if (CCW(p, q.front()) != CLOCKWISE) break;
		q.pop_back();
	}
	// 对称：用队尾线再清一遍队首
	while (SZ(q) >= 3) {
		auto p = getintersect(q.begin()[0], q.begin()[1]);
		if (CCW(p, q.back()) != CLOCKWISE) break;
		q.pop_front();
	}
	if (SZ(q) < 3) return {};
	vector<Point<DB>> poly;
	for (int i = 0; i < SZ(q); ++i) {
		poly.push_back(getintersect(q[i], q[(i + 1) % SZ(q)]));
	}
	return poly;
}
		\end{minted}

		\textbf{使用范例}（Uyuw's Concert 框架：给定 $N$ 条有向直线 + 题目自带的 $[0, 10000]^2$ 边界，求公共半平面交的面积）：

		\begin{minted}{cpp}
ll N; cin >> N;
vector<Line<DB>> lines(N);
for (auto &l : lines) cin >> l;
// 手动补 4 条边界直线（CCW 顺序：左下→右下→右上→左上）
auto a1 = Point<DB>(0, 0), a2 = Point<DB>(10000, 0);
auto a3 = Point<DB>(10000, 10000), a4 = Point<DB>(0, 10000);
lines.push_back(Line<DB>(a1, a2));
lines.push_back(Line<DB>(a2, a3));
lines.push_back(Line<DB>(a3, a4));
lines.push_back(Line<DB>(a4, a1));
auto poly = half_plane_cut(lines);
cout << fsp(1) << polygonArea(poly) << "\n";
		\end{minted}

		\subsection{凸包（Andrew 算法）}

		按 $x$ 升序、$x$ 相同按 $y$ 升序排序后，从左到右扫一遍维护下凸包，再从右到左扫一遍维护上凸包，把两段拼起来即可。

		\textbf{\texttt{strict} 参数}决定共线点处理：\texttt{strict = true} 时 \texttt{cross == 0}（共线）也弹栈，\textbf{不保留共线点}；\texttt{strict = false} 时只在 \texttt{cross < 0}（凹陷）才弹，\textbf{保留共线点}（题目要求「凸包上所有点」时用）。

		\textbf{\texttt{HullMode} 参数}：三种返回都是 CCW 方向。\texttt{HULL\_FULL} 完整凸包（\textbf{按 $Y$ 排序}，从\textbf{最低点}（$y$ 最小、再 $x$ 最小）起 CCW 绕一圈；这样剥层第一步就拿到最低点起步，不必再调 \texttt{reorder\_polygon} 找最低点）；\texttt{HULL\_LOWER} 下凸壳（最左 $\to$ 最右）；\texttt{HULL\_UPPER} 上凸壳（最右 $\to$ 最左，从已建好的全凸包数组里切片，不重跑单调栈）。\textbf{易踩}：\texttt{HULL\_UPPER} 等价「直线上半包络」时调用前需 dedupe 同斜率（strict 处理不了「同 $x$ 不同 $y$ 竖边」退化），见 STL 节 \texttt{sort + unique} 模板。

		\begin{minted}{cpp}
// 依赖：
//   struct Point
//   cross(Point, Point) 与三参数 cross(Point, Point, Point)
enum HullMode { HULL_FULL, HULL_LOWER, HULL_UPPER };

// strict = true:  严格凸包 (不保留三点共线的中间点)
// strict = false: 弱凸包    (保留三点共线的中间点), 题目要求"凸包上所有点"时用
template<class T>
vector<Point<T> > make_convex_hull(vector<Point<T> > A, bool strict, HullMode hull_mode) {
	// Andrew 扫描. 各 mode 的排序键与返回 (顶点序列都是 CCW 序):
	// HULL_FULL : 按 Y 排序 (Y 同按 X), 返回整圈凸包, 从 (y 最小, 再 x 最小) 点起头.
	// HULL_LOWER: 按 X 排序 (X 同按 Y), 返回下凸链, 最左 -> 最右 (x 单调不减).
	// HULL_UPPER: 按 X 排序 (X 同按 Y), 返回上凸链, 最右 -> 最左 (x 单调不增).
	if (hull_mode == HULL_FULL) {
		sort(A.begin(), A.end(), [](const Point<T> &a, const Point<T> &b) {
			return a.y != b.y ? a.y < b.y : a.x < b.x;
		});
	} else {
		sort(A.begin(), A.end());
	}
	A.resize(unique(all(A)) - A.begin());
	size_t n = A.size();
	if (n <= 2) {
		// 退化: 不足 3 点直接返回排序后的点; UPPER 逆序成 x 不增, FULL/LOWER 保持排序序
		if (hull_mode == HULL_UPPER) reverse(A.begin(), A.end());
		return A;
	}
	vector<Point<T> > ans(n * 2);
	// strict 时 cross == 0 共线也弹; 弱凸只在 cross < 0 凹陷时弹
	auto fail = [&](ll s) { return strict ? s >= 0 : s > 0; };
	int now = -1;
	for (int i = 0; i < n; i++) {
		// 维护下凸包
		while (now > 0 && fail(sgn(cross(A[i], ans[now], ans[now - 1])))) {
			now--; // 弹出栈顶，因为它会导致凹陷 (strict 时还会弹掉共线)
		}
		ans[++now] = A[i];
	}
	if (hull_mode == HULL_LOWER) {
		ans.resize(now + 1);
		return ans;
	}
	int pre = now;
	// 从右向左扫描所有点（注意是从 n-2 开始，因为 n-1 已经在栈顶了）
	for (int i = n - 2; i >= 0; i--) {
		// 维护上凸包
		while (now > pre && fail(sgn(cross(A[i], ans[now], ans[now - 1])))) {
			now--;
		}
		ans[++now] = A[i];
	}
	if (hull_mode == HULL_UPPER) {
		return vector<Point<T> >(ans.begin() + pre, ans.begin() + now + 1);
	}
	ans.resize(now);
	// HULL_FULL: 因上面按 Y 排序, 这里已天然是「Y 最小点起头 + CCW」.
	// 退化处理: 输入点全共线时凸包退化成线段, 弱凸包会折返成重复序列 (如 A B C B);
	// 有向面积为 0 即退化, 此时 A 已按 Y 排好序且点互异, 就是沿线单程, 直接返回 (O(n), 不折返).
	if (sgn(polygonSignedArea(ans)) == 0 && !strict) {
		return A;
	}
	return ans;
}
		\end{minted}

		\subsubsection{对偶点技巧：直线上半包络 $\Leftrightarrow$ 对偶点上凸壳}

		\href{https://www.luogu.com.cn/problem/P3194}{洛谷 P3194 [HNOI2008] 水平可见直线}：给 $N$ 条直线 $y = A_i x + B_i$，从 $y = +\infty$ 俯视，问哪些「有某段 $x$ 露头」。

		\textcolor{blue}{\textbf{直线-点对偶}}：把 $L_i$ 映射到对偶平面的点 $P_i = (A_i, B_i)$，则
		\[
			L_i \text{ 可见} \iff P_i \text{ 在所有对偶点的「严格上凸壳」上.}
		\]
		\textbf{为什么}：3 条 $a_1 < a_2 < a_3$，$L_2$ 被 $L_1, L_3$ 在某 $x$ 处压下 $\iff \mathrm{cross}(P_2 - P_1,\, P_3 - P_1) \ge 0$（$P_2$ 在 $P_1 P_3$ 下方或共线）。共线时三线共点、$L_2$ 只在一点露头不是 segment，\texttt{strict} 模式正好弹掉。

		\begin{center}
			\begin{tikzpicture}[>=Stealth, font=\scriptsize]
				% ===== Primal plane =====
				\begin{scope}[x=0.55cm, y=0.45cm]
					\draw[->, gray] (-1.7, 0) -- (1.9, 0) node[right] {$x$};
					\draw[->, gray] (0, -1.3) -- (0, 3.6) node[left] {$y$};
					% 4 条线全画（细灰）
					\draw[blue!25, thin] (-1.5,  4)   -- (1.5, -2);   % L1: y = -2x + 1
					\draw[blue!25, thin] (-1.5,  2)   -- (1.5,  2);   % L2: y = 2
					\draw[blue!25, thin] (-1.5, -2)   -- (1.5,  4);   % L4: y = 2x + 1
					% 上半包络（粗蓝）：L1 → 转折(-0.5, 2) → L2 → 转折(0.5, 2) → L4
					\draw[blue, very thick] (-1.5, 4) -- (-0.5, 2) -- (0.5, 2) -- (1.5, 4);
					% L3 被压（灰虚）
					\draw[gray!70, dashed, thin] (-1.5, 2.5) -- (1.5, -0.5);
					% 标签
					\node[blue!70!black, font=\tiny] at (-1.55, 3.55) {$L_1$};
					\node[blue!70!black, font=\tiny] at ( 1.55, 3.55) {$L_4$};
					\node[blue!70!black, font=\tiny] at ( 1.55, 2.35) {$L_2$};
					\node[gray, font=\tiny]          at ( 1.65, -0.7) {$L_3$};
					\node[font=\footnotesize, anchor=south] at (0, 3.8) {\textbf{Primal}};
				\end{scope}

				% ===== Dual plane =====
				\begin{scope}[xshift=4.6cm, x=0.55cm, y=0.7cm]
					\draw[->, gray] (-2.7, 0) -- (2.9, 0) node[right] {$A$};
					\draw[->, gray] (0, -0.4) -- (0, 2.7) node[left] {$B$};
					% 上凸壳（粗蓝）
					\draw[blue, very thick] (-2, 1) -- (0, 2) -- (2, 1);
					% 提示 P3 严格在弦下方的垂直差
					\draw[red!60, dashed, thin] (-1, 1) -- (-1, 1.5);
					% 4 个点
					\fill[blue] (-2, 1) circle (2pt) node[below left,  font=\tiny, color=blue] {$P_1$};
					\fill[blue] ( 0, 2) circle (2pt) node[above,       font=\tiny, color=blue] {$P_2$};
					\fill[blue] ( 2, 1) circle (2pt) node[below right, font=\tiny, color=blue] {$P_4$};
					% P3 被弹（红叉）
					\fill[gray] (-1, 1) circle (1.8pt);
					\draw[red, thick] (-1.18, 0.82) -- (-0.82, 1.18);
					\draw[red, thick] (-1.18, 1.18) -- (-0.82, 0.82);
					\node[gray, font=\tiny, anchor=north] at (-1, 0.85) {$P_3$};
					\node[font=\footnotesize, anchor=south] at (0, 2.85) {\textbf{Dual}};
				\end{scope}
			\end{tikzpicture}

			\vspace{2pt}
			{\footnotesize\sffamily\color{gray} 图注：左 Primal 4 条线，上半包络 = $L_1 \to L_2 \to L_4$（粗蓝折线），$L_3$ 全程在下被压（灰虚线）。右 Dual 4 个对偶点，$P_3 = (-1, 1)$ 严格落在弦 $P_1 P_2$ 下方（红虚示意垂直差 0.5），上凸壳 $P_1 \to P_2 \to P_4$ 不收 $P_3$。\textbf{这就是「$L_i$ 可见 $\Leftrightarrow$ $P_i$ 在上凸壳上」的视觉对照}。}
		\end{center}

		\textbf{落地}：(1) 把 $(A_i, B_i)$ 当点（套 \texttt{Point<ll>} 的 \texttt{id} 字段绑原始下标）；(2) 同 $A$ 多条只留 $B$ 最大那条（\texttt{sort + unique} 模板见 STL 节）；(3) \texttt{make\_convex\_hull(pts, true, HULL\_UPPER)} 直接返回上凸壳点列，按 \texttt{id} 升序输出。$|A|, |B| \le 5 \times 10^5$ 时全程整数，\texttt{cross} 走 \texttt{\_\_int128}。

		\subsection{凸多边形直径（旋转卡壳）}

		\href{https://vjudge.net/problem/QOJ-784}{QOJ 784}（凸多边形直径模板题）。

		给一个 \textbf{逆时针顺序}的 $n$ 凸多边形（允许三点共线），求最远点对的距离。\textcolor{blue}{\textbf{旋转卡壳}}的标准技巧：把每条边 $A[r] \to A[r{+}1]$ 当成一条「卡尺」，对踵点 $A[l]$ 是离这条边最远的顶点。$r$ 沿 CCW 扫一圈时 $l$ 单调（不回退），整体 $O(n)$。

		\textbf{关键判据}：$l$ 该不该往前推一步？看「以边 $(r,\,r{+}1)$ 为底、$l{+}1$ vs $l$ 为顶」的两个三角形面积——后者 $\ge$ 前者就推。\textbf{比直接比距离平方更鲁棒}：曲率非均匀的凸包（如 spike 形 $r(\theta) = 1 + 0.1\sin(7\theta)$）下 $|A[l] - A[r]|^2$ 沿 $r$ 不是单峰会卡死，但「以边为底的三角形面积」仍是单峰。

		\textbf{4 个不要踩的坑}（QOJ 784 session 实录，每个都有真实反例）：
		\begin{enumerate}
					\setlength\itemsep{0pt}
				\item \textbf{不要做点对踵}：以 $|A[l] - A[r]|^2$ 沿 $r$ 找单峰只在「曲率均匀」凸包成立，cyaron 默认凸包形是盲区——必死于 spike 形；
				\item \textbf{$l, r$ 角色不能反}：必须「外层 $r$ 是边、内层 $l$ 是点」，反过来 $|\text{edge}|$ 乘子会破坏单峰；
				\item \textbf{ans 候选取两个}：$(A[l], A[r])$ 和 $(A[l], A[r{+}1])$ 都要更新，只取一个会漏掉「边端点本身就是最远对」的情形；
				\item \textbf{推进用 $\le$ 不是 $<$}：plateau（两个候选三角形面积相等）起点要继续推到 plateau 末端，严格 $<$ 会卡在起点。
		\end{enumerate}

		\textbf{对拍提醒}：几何题对拍\textbf{必须用曲率非均匀的形状}（spike / heart / lemon / 椭圆），cyaron \texttt{Polygon.convex\_hull} 默认是近圆形，stress 5000+ case 全过会让你错误推断「算法对」。详见 \texttt{cyaron-stress} skill。

		\begin{minted}{cpp}
// 依赖：
//   struct Point, sqrtL（基础模板已定义）
//   cross（含三参数）, isConvex（LOCAL 校验用，弱凸口径）
//   宏 SZ, all（用户自定义）

// 距离平方（不开方）：旋转卡壳里 ans 反复更新，留到最后一次 sqrt
template<class T, class U>
auto distancePPNorm(const Point<T> &a, const Point<U> &b) {
	return (a - b).norm();
}

// L2 欧几里得距离 (sqrt 版, 用 long double). 整型坐标走 isqrt + fraction 不丢精度.
template<class T, class U>
DB distancePP(const Point<T> &a, const Point<U> &b) {
	return (a - b).abs();
}

// L1 曼哈顿距离 = |dx| + |dy|. 返回类型继承 normL1 (整型 -> __int128_t, 浮点 -> T).
template<class T, class U>
auto distancePPL1(const Point<T> &a, const Point<U> &b) {
	return (a - b).normL1();
}

// Linf 切比雪夫距离 = max(|dx|, |dy|). 返回类型 T.
template<class T, class U>
T distancePPLinf(const Point<T> &a, const Point<U> &b) {
	return (a - b).normLinf();
}

// 以 CCW 顺序给定 n 点凸多边形（允许三点共线），求直径（最远点对距离）
template<class T>
long double convex_diameter(const vector<Point<T>> &A) {
#ifdef LOCAL
	vector<Point<T>> B(all(A));
	assert(isConvex(B, false)); // 弱凸即可，允许三点共线
#endif
	// 退化输入特判：n <= 2 时 area() 恒 0 会让 need_l_move 永真死循环
	if (SZ(A) <= 1) return 0;
	if (SZ(A) == 2) return sqrtL(distancePPNorm(A[0], A[1]));
	auto ans = (A[0] - A[0]).norm(); // 零值，类型由 norm 继承
	// 以边 (r, r+1) 为底，A[l] 为顶的三角形 2 倍面积
	auto area = [&](ll l, ll r) {
		auto res = cross(A[r], A[(r + 1) % SZ(A)], A[l]);
		if (res < 0) res = -res;
		return res;
	};
	// l 是否还要向前推？比较 (l, r) vs (l+1, r) 三角形面积
	// 用 <= 推到 plateau 末端，严格 < 会卡在 plateau 起点
	auto need_l_move = [&](ll l, ll r) -> bool {
		return area(l, r) <= area((l + 1) % SZ(A), r);
	};
	for (int l = 0, r = 0; r < SZ(A); ++r) {
		while (need_l_move(l, r)) l = (l + 1) % SZ(A);
		ans = max(ans, distancePPNorm(A[l], A[r]));
		ans = max(ans, distancePPNorm(A[l], A[(r + 1) % SZ(A)])); // 第二候选
	}
	return sqrtL(ans);
}
		\end{minted}

		\subsection{最远点对（旋转卡壳，返回原始下标）}

		\href{https://www.luogu.com.cn/problem/P16223}{Luogu P16223}。

		\texttt{convex\_diameter} 只返回距离；要把最远点对的原始输入下标也输出回来，靠 \texttt{Point::id} 字段：读入时填 \texttt{A[i].id = i}，\texttt{make\_convex\_hull} 不做派生运算（仅 \texttt{sort} + 拷贝原点入栈），凸包顶点继承原始 \texttt{id}。

		\begin{minted}{cpp}
// 依赖：
//   struct Point（带 id 字段）
//   cross（三参数）, distancePPNorm, make_convex_hull
//   宏 SZ（用户自定义）
template<class T>
pair<int, int> furthest_pair_of_point(vector<Point<T>> A) {
	A = make_convex_hull(A, true, HULL_FULL);
	// 退化短路：hull == 1 全重合；hull == 2 全共线（area=0 死循环）
	if (SZ(A) == 1) return {A[0].id, A[0].id};
	if (SZ(A) == 2) return {A[0].id, A[1].id};
	i128 ans = 0;
	auto area = [&](ll l, ll r) {
		auto res = cross(A[r], A[(r + 1) % SZ(A)], A[l]);
		if (res < 0) res = -res;
		return res;
	};
	auto need_l_move = [&](ll l, ll r) -> bool {
		return area(l, r) <= area((l + 1) % SZ(A), r);
	};
	pair<int, int> ans_pair;
	for (int l = 0, r = 0; r < SZ(A); ++r) {
		while (need_l_move(l, r)) l = (l + 1) % SZ(A);
		if (distancePPNorm(A[l], A[r]) > ans) {
			ans = distancePPNorm(A[l], A[r]);
			ans_pair = {A[l].id, A[r].id};
		}
		int rn = (r + 1) % SZ(A);
		if (distancePPNorm(A[l], A[rn]) > ans) {
			ans = distancePPNorm(A[l], A[rn]);
			ans_pair = {A[l].id, A[rn].id};
		}
	}
	return ans_pair;
}
		\end{minted}

		\textbf{调用样板}（读入时给 \texttt{id}）：
		\begin{minted}{cpp}
ll N; cin >> N;
vector<Point<ll>> A(N);
for (int i = 0; i < N; ++i) {
	cin >> A[i];
	A[i].id = i;
}
auto [l, r] = furthest_pair_of_point(A);
cout << l << " " << r << '\n';
		\end{minted}

		\subsection{最小矩形覆盖（旋转卡壳，4 指针）}

		\href{https://vjudge.net/problem/Luogu-P3187}{Luogu P3187 [HNOI2007] 最小矩形覆盖}。

		给一个点集，求面积最小、把所有点框住的矩形（不必轴对齐）。\textbf{关键定理}：最小覆盖矩形必有一条边与凸包某条边重合。所以枚举每条凸包边 $e$ 当矩形底边，三指针单调推进维护：
		\begin{itemize}
				\setlength\itemsep{0pt}
			\item \texttt{up}：以 $e$ 为底时离 $e$ 最远的顶点（决定矩形高 $h$）；
			\item \texttt{r} / \texttt{l}：在 $e$ 方向上投影最大 / 最小的顶点（决定宽 $w$ 的右 / 左端）。
		\end{itemize}
		外层 \texttt{edge} 沿 CCW 扫一圈，内层 \texttt{up}, \texttt{r}, \texttt{l} 各自单调（不回退），整体 $O(n)$。

		\textbf{比较时不要早除}：实现里用 \texttt{w = dot(e, A[r] - A[l])}（长度乘子 $|e|$）和 \texttt{h = cross(e, A[up] - A[edge])}（同样乘子 $|e|$），所以 \texttt{area = w * h} 是真实矩形面积的 $|e|^2$ 倍。比较两条不同底边 $e$ 给出的候选时，正确做法是\textbf{交叉相乘}：
		\[
			\texttt{area} \cdot |e_{\text{best}}|^2 \quad\text{vs.}\quad \texttt{mn\_area} \cdot |e|^2
		\]
		这样比较时只走整数（\texttt{T = DB} 下也是平方量级），避免提前除 $|e|^2$ 丢精度。\textbf{还原 4 顶点}：底 / 顶切线方向 $\parallel e$，左 / 右切线方向 $\perp e$，4 条切线两两相交即得。函数签名\textbf{硬钉 \texttt{Point<DB>}}（\texttt{getintersect} 走浮点除法 + 还原 4 顶点必然是非整坐标，没有保留整数 \texttt{T} 的意义；类型层一刀切防"误传 \texttt{vector<Point<ll>>}"沉默截断）。

		\begin{minted}{cpp}
// 依赖：
//   struct Point, struct Line
//   cross（含三参数）, dot, getintersect
//   make_convex_hull(strict = true), Point::norm()
//   宏 SZ（用户自定义）
// 最小矩形覆盖：返回 {面积, 4 个矩形顶点（以底边 BL 起、CCW）}
// 前置：T = DB（reconstruct 走 getintersect 必须浮点）
pair<DB, array<Point<DB>, 4> > minimum_rectangle_cover(vector<Point<DB> > A) {
	A = make_convex_hull(A, true, HULL_FULL);
	int n = SZ(A);
	auto e_of = [&](int i) { return A[(i + 1) % n] - A[i]; }; // 第 i 条凸包边的方向向量
	int up = 0, l = 0, r = 0;
	using Type = decltype(dot(A[0], A[0]));
	Type mn_area = 0;
	do {
		auto e0 = e_of(0);
		for (int i = 0; i < SZ(A); ++i) {
			// 同一个底，面积越大，高自然越大
			if (cross(e0, A[i] - A[0]) > cross(e0, A[up] - A[0])) {
				up = i;
			}
			if (dot(e0, A[i] - A[0]) > dot(e0, A[r] - A[0])) {
				r = i;
			}
			if (dot(e0, A[i] - A[0]) < dot(e0, A[l] - A[0])) {
				l = i;
			}
		}
		auto w = dot(e0, A[r] - A[l]);
		auto h = cross(e0, A[up] - A[0]);
		mn_area = w * h;
	} while (false);
	int bestL = l, bestR = r, bestE = 0, bestUp = up;
	for (int edge = 0; edge < SZ(A); ++edge) {
		Point<DB> e = e_of(edge);
		while (cross(e, A[(up + 1) % SZ(A)] - A[edge]) > cross(e, A[up] - A[edge]))
			up = (up + 1) % SZ(A);
		while (dot(e, A[(r + 1) % SZ(A)] - A[edge]) > dot(e, A[r] - A[edge]))
			r = (r + 1) % SZ(A);
		while (dot(e, A[(l + 1) % SZ(A)] - A[edge]) < dot(e, A[l] - A[edge]))
			l = (l + 1) % SZ(A);
		// 这个宽和高均乘以了 e 的长度
		auto w = dot(e, A[r] - A[l]);
		auto h = cross(e, A[up] - A[edge]);
		auto area = w * h;
		if (area * e_of(bestE).norm() < mn_area * e.norm()) {
			mn_area = area;
			// mn_area_edge = e.norm();
			tie(bestL, bestR, bestE, bestUp) = make_tuple(l, r, edge, up);
		}
	}
	// 还原 4 顶点 = 4 条切线两两相交（4 条线方向：底/顶 parallel e；左/右 parallel e_perp）
	Point e = e_of(bestE);
	Point e_perp(-e.y, e.x); // CCW 90° 旋转
	Line bottom(A[bestE], A[bestE] + e);
	Line top_(A[bestUp], A[bestUp] + e);
	Line right_(A[bestR], A[bestR] + e_perp);
	Line left_(A[bestL], A[bestL] + e_perp);
	array<Point<DB>, 4> corners = {
		getintersect(bottom, left_), // BL
		getintersect(bottom, right_), // BR
		getintersect(top_, right_), // TR
		getintersect(top_, left_), // TL
	};
	return {(long double) mn_area / e.norm(), corners};
}
		\end{minted}

		\subsection{平面最近点对}

		经典分治：把所有点按 $x$ 排序后，递归求左右两半的最小距离 $d$，再处理跨越中线、与中线 $x$ 距离 $< \sqrt{d}$ 的「窄带」候选点。

		\textbf{难点}在于把窄带内的候选点收集起来后再按 $y$ 暴力扫一遍——技巧是从中线向两侧分别推进，一旦 $|\Delta x| \ge \sqrt{d}$ 立刻 \texttt{break}，避免把窄带退化成 $O(n)$。

		\begin{minted}{cpp}
// 依赖：
//   struct Point（基础模板，本节用 Point<ll>）
//   dot(Point, Point)
//   宏 SZ, FOR, ROF, EB（用户自定义）
//   全局数组 Point<ll> A[MAXN]
Point<ll> A[MAXN];
struct Best {
	i128 d;        // dot<ll>(v, v) 返回 i128，得收得下，否则 1e9 坐标平方就溢出 ll
	Point<ll> a, b;
	bool operator<(const Best &o) const { return d < o.d; }
};

// 暴力求平面最近点对（按 y 排序后扫）
Best mindis(vector<Point<ll>> &a, Best res) {
	sort(all(a), [&](Point<ll> a, Point<ll> b) { return a.y < b.y; });
	FOR(i, 1, SZ(a) - 1) {
		ROF(j, i - 1, 0) {
			Point<ll> p = a[i], q = a[j];
			// y 距离平方已超出当前最优就早停
			if ((p.y - q.y) * (p.y - q.y) > res.d) break;
			Point<ll> v = p - q;
			res = min(res, {dot(v, v), p, q});
		}
	}
	return res;
}

Best dfs(int l, int r) {
	if (l >= r) return {(i128) INF, A[l], A[l]};
	if (r - l == 1) {
		Point<ll> v = A[r] - A[l];
		return {dot(v, v), A[l], A[r]};
	}
	int mid = r + l >> 1;
	Best ret = dfs(l, mid);
	ret = min(ret, dfs(mid + 1, r));
	ll x = A[mid + 1].x;
	vector<Point<ll>> wp;
	// 左半向左推进，超出 sqrt(ret.d) 就停
	ROF(i, mid, l) {
		i128 d = abs(A[i].x - x);
		if (d * d < ret.d) wp.EB(A[i]);
		else break;
	}
	// 右半向右推进
	FOR(i, mid + 1, r) {
		i128 d = abs(A[i].x - x);
		if (d * d < ret.d) wp.EB(A[i]);
		else break;
	}
	return mindis(wp, ret);
}
		\end{minted}

		\subsection{两凸多边形最近距离（旋转卡壳）}

		\href{https://vjudge.net/problem/OpenJ_Bailian-3851}{OpenJ\_Bailian 3851 Bridge Across Islands}。给两个不相交的 CCW 凸多边形 $A, B$，求两边界最小距离。

		初始 $\texttt{eA}$ 取 $A$ 最低点、$\texttt{eB}$ 取 $B$ 最高点（一对反向支撑线起步）；按 $\texttt{cross}(\vec{e_A}, \vec{e_B}) > 0$ 单调推进 $\texttt{eB}$、否则推进 $\texttt{eA}$，合计 $\le n + m$ 步。每步取段-段距离 \texttt{distanceSS}（A 边长且 B 边整段在 A 边投影内时单纯点-边距离会漏，必须段-段）。

		\textcolor{red}{\textbf{必须双向各调一次}}：单次 \texttt{cal\_min\_dis\_two\_convex(A, B)} 的对偶覆盖不完整，要 \texttt{(A, B)} 和 \texttt{(B, A)} 各跑一次取 $\min$ 才覆盖全。漏一次会在「$A$ 顶点 → $B$ 边」和「$B$ 顶点 → $A$ 边」非对称分布下系统性偏大（实测反例存在）。

		\begin{minted}{cpp}
// AC https://vjudge.net/solution/69793877
// 依赖：
//   struct Point, struct Line
//   reorder_polygon, cross, distanceSS
//   宏 SZ, all（用户自定义）
template<class T>
DB cal_min_dis_two_convex(vector<Point<T> > A, vector<Point<T> > B) {
	reorder_polygon(A);
	reorder_polygon(B);
	ll yMnA = min_element(all(A), [](const Point<T> &a, const Point<T> &b) {
		return a.y < b.y;
	}) - A.begin(), yMxB = max_element(all(B), [](const Point<T> &a, const Point<T> &b) {
		return a.y < b.y;
	}) - B.begin();
	DB ans = INFINITY;
	ll eA = yMnA, eB = yMxB;
	for (int _ = 0; _ < SZ(A) + SZ(B); ++_) {
		auto check = [&]() {
			auto vecA = A[(eA + 1) % SZ(A)] - A[eA];
			auto vecB = B[(eB + 1) % SZ(B)] - B[eB];
			return cross(vecA, vecB) > 0;
		};
		while (check()) ++eB, eB %= SZ(B);
		ans = min(ans, distanceSS(Line(A[(eA + 1) % SZ(A)], A[eA]),
								  Line(B[(eB + 1) % SZ(B)], B[eB])));
		eA += 1;
		eA %= SZ(A);
	}
	return ans;
}

// 调用方：必须双向各跑一次取 min
// DB ans = min(cal_min_dis_two_convex(A, B), cal_min_dis_two_convex(B, A));
		\end{minted}

	\newpage

	\subsection{常用公式}

	\subsubsection{正多边形面积公式}

	设边数为 $n$、边长为 $l$：
	$$
	A = \dfrac{n \cdot l^2}{4 \cdot \tan\!\left(\dfrac{\pi}{n}\right)}
	$$

	\subsubsection{三角恒等式}

	\textbf{和差角}：
	\begin{align*}
		\sin(\alpha \pm \beta) & = \sin\alpha\cos\beta \pm \cos\alpha\sin\beta                                  \\
		\cos(\alpha \pm \beta) & = \cos\alpha\cos\beta \mp \sin\alpha\sin\beta                                  \\
		\tan(\alpha \pm \beta) & = \dfrac{\tan\alpha \pm \tan\beta}{1 \mp \tan\alpha\tan\beta}
	\end{align*}

	\textbf{倍角}：
	$$
	\sin 2\theta = 2\sin\theta\cos\theta, \qquad
	\cos 2\theta = \cos^2\theta - \sin^2\theta = 1 - 2\sin^2\theta = 2\cos^2\theta - 1
	$$
	$$
	\tan 2\theta = \dfrac{2\tan\theta}{1 - \tan^2\theta}
	$$

	\textbf{降幂（把 $\theta$ 整体升到 $2\theta$）}：
	$$
	\sin^2\theta = \dfrac{1 - \cos 2\theta}{2}, \qquad
	\cos^2\theta = \dfrac{1 + \cos 2\theta}{2}, \qquad
	\sin\theta\cos\theta = \dfrac{\sin 2\theta}{2}
	$$

	\textbf{辅助角（线性叠加 $\to$ 单一正弦）}：
	$$
	a\sin\theta + b\cos\theta = \sqrt{a^2 + b^2}\,\sin(\theta + \varphi),
	\qquad \tan\varphi = \dfrac{b}{a}
	$$
	极值即 $\pm\sqrt{a^2 + b^2}$，把单变量旋转优化变成闭式。

	\textbf{积化和差}：
	\begin{align*}
		\sin\alpha\cos\beta & = \tfrac{1}{2}\bigl[\sin(\alpha+\beta) + \sin(\alpha-\beta)\bigr] \\
		\cos\alpha\cos\beta & = \tfrac{1}{2}\bigl[\cos(\alpha+\beta) + \cos(\alpha-\beta)\bigr] \\
		\sin\alpha\sin\beta & = \tfrac{1}{2}\bigl[\cos(\alpha-\beta) - \cos(\alpha+\beta)\bigr]
	\end{align*}

	\textbf{和差化积}：
	\begin{align*}
		\sin\alpha + \sin\beta & = 2\sin\tfrac{\alpha+\beta}{2}\cos\tfrac{\alpha-\beta}{2}  \\
		\sin\alpha - \sin\beta & = 2\cos\tfrac{\alpha+\beta}{2}\sin\tfrac{\alpha-\beta}{2}  \\
		\cos\alpha + \cos\beta & = 2\cos\tfrac{\alpha+\beta}{2}\cos\tfrac{\alpha-\beta}{2}  \\
		\cos\alpha - \cos\beta & = -2\sin\tfrac{\alpha+\beta}{2}\sin\tfrac{\alpha-\beta}{2}
	\end{align*}

	\textbf{半角}：
	$$
	\sin\dfrac{\theta}{2} = \pm\sqrt{\dfrac{1-\cos\theta}{2}}, \qquad
	\cos\dfrac{\theta}{2} = \pm\sqrt{\dfrac{1+\cos\theta}{2}}, \qquad
	\tan\dfrac{\theta}{2} = \dfrac{\sin\theta}{1+\cos\theta} = \dfrac{1-\cos\theta}{\sin\theta}
	$$

	\newpage

	\subsection{附录：C++14 版本基础模板}

	部分老评测机只开到 C++14（POJ / HDU / UOJ 早期等），\texttt{if constexpr}、变量模板别名 \texttt{is\_xxx\_v}、\texttt{inline constexpr} 变量都不能用。下面给出与正文「点线基础模板」语义等价的 C++14 重写，赛场上换评测机直接整段替换基础模板即可，后续 \texttt{make\_convex\_hull} / \texttt{minkowski} / 半平面交 / \texttt{Circle} 等子节代码\textbf{无需改动}。

	四点改造原则：\textbf{(1)} \texttt{std::is\_xxx\_v} 是 C++17 才有的别名，手写一份放\textbf{全局命名空间}（不污染 \texttt{std}）—— C++14 的 \texttt{std} 里没有这层别名，所以 \texttt{using namespace std;} 之后下游代码可以原封不动写 \texttt{is\_floating\_point\_v<T>}。\textbf{(2)} \texttt{sqrtL} / \texttt{squared} / \texttt{cross} / \texttt{dot} 等本来靠 \texttt{if constexpr} 在函数体内做整型/浮点分流的函数模板，全部拆成两个独立模板，用 \texttt{enable\_if\_t<...>} 写互补条件覆盖所有 \texttt{T} 组合。\textbf{(3)} 成员函数 \texttt{Point::norm} / \texttt{Point::normL1} 不能直接用 class 的 \texttt{T} 做 SFINAE（class 实例化时就会 hard error），借 \texttt{template<class TT = T, enable\_if\_t<...>>} 把 SFINAE 推迟到成员函数自身的模板实例化阶段。\textbf{(4)} \texttt{is\_promotion\_v} 去掉 \texttt{inline} —— C++14 已有变量模板，但变量不能 \texttt{inline}（单 TU 内无 ODR 冲突，多 TU 才需要 C++17 \texttt{inline}）。

	\begin{minted}{cpp}
// === C++14 版本基础模板 (POJ / HDU / UOJ 等老评测机) ===
// 1) std::xxx_v 是 C++17 才有的别名, 这里放全局命名空间手写一份;
//    C++14 std 里没这层别名, 与 using namespace std 不冲突
template<class T> constexpr bool is_floating_point_v = std::is_floating_point<T>::value;
template<class T> constexpr bool is_arithmetic_v     = std::is_arithmetic<T>::value;
template<class T> constexpr bool is_integral_v       = std::is_integral<T>::value;
template<class F, class T> constexpr bool is_same_v  = std::is_same<F, T>::value;

// 2) isqrt: 不依赖 C++17, 与正文版完全一致
i128 isqrt(i128 n) {
	if (n <= 0) return 0;
	i128 lo = 0, hi = (i128) 2e18;
	while (lo < hi) {
		i128 mid = lo + (hi - lo + 1) / 2;
		if (mid <= n / mid) lo = mid;
		else                hi = mid - 1;
	}
	return lo;
}

// 3) sqrtL: if constexpr 拆两个独立 SFINAE 模板, 互补条件覆盖所有 T
template<typename T, enable_if_t<is_floating_point_v<T>, int> = 0>
long double sqrtL(T n) { return sqrtl((long double) n); }

template<typename T, enable_if_t<!is_floating_point_v<T>, int> = 0>
long double sqrtL(T n) {
	i128 ni = (i128) n;
	if (ni <= 0) return 0;
	i128 i = isqrt(ni);
	long double id = (long double) i;
	long double rd = (long double) (ni - i * i);
	return id * sqrtl(1.0L + rd / (id * id));
}

// 4) squared: 同 sqrtL 拆法
template<class T, enable_if_t<is_floating_point_v<T>, int> = 0>
auto squared(T x) { return x * x; }

template<class T, enable_if_t<!is_floating_point_v<T>, int> = 0>
i128 squared(T x) { return (i128) x * (i128) x; }

// sgn: 不依赖 C++17, 照搬
template<class T>
int sgn(T x) {
	if (-eps < x && x < eps) return 0;
	return x < 0 ? -1 : 1;
}

// 5) Point: 成员函数 norm / normL1 借 TT = T 默认模板参数,
//    把 SFINAE 推迟到成员函数实例化, 避免 class 实例化时直接 hard error
//    converting ctor / 标量乘除走 is_same_v<common_type_t<U,T>, T> 方向限制
template<class T>
struct Point {
	T x, y;
	int id = -1;

	explicit Point(T x_ = 0, T y_ = 0) : x(x_), y(y_) {}

	// Point<U> -> Point<T> 隐式 (common_type 单调升档方向, U -> T 必须无损):
	// common_type_t<U, T> == T 才放行 (例 ll -> DB 通; DB -> ll 拦, 防 truncation)
	template<class U, class = enable_if_t<is_same_v<common_type_t<U, T>, T>>>
	Point(const Point<U> &o) : x((T) o.x), y((T) o.y), id(o.id) {}

	Point &operator+=(Point p) & { x += p.x, y += p.y; return *this; }
	Point &operator-=(Point p) & { x -= p.x, y -= p.y; return *this; }
	Point &operator*=(T t) & { x *= t, y *= t; return *this; }
	Point &operator/=(T t) & { x /= t, y /= t; return *this; }

	Point operator-() const { return Point(-x, -y); }

	friend Point operator+(Point a, Point b) { return a += b; }
	friend Point operator-(Point a, Point b) { return a -= b; }

	// 标量乘除: 返回 Point<common_type_t<T, U>>; R==T 走 copy ctor, R!=T 走 converting ctor
	template<class U>
	friend auto operator*(Point a, U b) {
		Point<common_type_t<T, U>> r = a; r *= b; return r;
	}
	template<class U>
	friend auto operator*(U a, Point b) {
		Point<common_type_t<T, U>> r = b; r *= a; return r;
	}
	template<class U>
	friend auto operator/(Point a, U b) {
		Point<common_type_t<T, U>> r = a; r /= b; return r;
	}

	friend bool operator<(Point a, Point b) {
		int sx = sgn(a.x - b.x);
		if (sx != 0) return a.x < b.x;
		return a.y < b.y;
	}
	friend bool operator>(Point a, Point b) { return b < a; }
	friend bool operator==(Point a, Point b) {
		return sgn(a.x - b.x) == 0 && sgn(a.y - b.y) == 0;
	}
	friend bool operator!=(Point a, Point b) { return !(a == b); }
	friend istream &operator>>(istream &is, Point &p) { return is >> p.x >> p.y; }

	// 关键: TT = T 让 SFINAE 推迟到成员函数实例化阶段
	template<class TT = T, enable_if_t<!is_floating_point_v<TT>, int> = 0>
	auto norm() const { return (__int128_t) x * x + (__int128_t) y * y; }
	template<class TT = T, enable_if_t<is_floating_point_v<TT>, int> = 0>
	auto norm() const { return x * x + y * y; }

	DB abs() const { return sqrtL(norm()); }

	Point rot90() const { return Point(-y, x); }
	Point<DB> rot(DB ang) const {
		return Point<DB>(cosl(ang) * x - sinl(ang) * y,
		                 sinl(ang) * x + cosl(ang) * y);
	}

	template<class TT = T, enable_if_t<!is_floating_point_v<TT>, int> = 0>
	auto normL1() const {
		auto ax = (x < 0 ? -x : x), ay = (y < 0 ? -y : y);
		return (__int128_t) ax + (__int128_t) ay;
	}
	template<class TT = T, enable_if_t<is_floating_point_v<TT>, int> = 0>
	auto normL1() const {
		auto ax = (x < 0 ? -x : x), ay = (y < 0 ? -y : y);
		return ax + ay;
	}

	T normLinf() const {
		auto ax = (x < 0 ? -x : x), ay = (y < 0 ? -y : y);
		return ax > ay ? ax : ay;
	}
};

// 7) cross (两参 / 三参): 两参拆两个 SFINAE 模板; 三参不需分流照搬
template<class T, enable_if_t<!is_floating_point_v<T>, int> = 0>
__int128_t cross(Point<T> a, Point<T> b) {
	return (__int128_t) a.x * b.y - (__int128_t) a.y * b.x;
}
template<class T, enable_if_t<is_floating_point_v<T>, int> = 0>
auto cross(Point<T> a, Point<T> b) {
	return a.x * b.y - a.y * b.x;
}

template<class T>
auto cross(Point<T> o, Point<T> a, Point<T> b) { return cross(a - o, b - o); }

// 8) dot: 与 cross 同构, 拆两个 SFINAE 模板 (T / U 可不同型)
template<class T, class U,
	enable_if_t<!is_floating_point_v<T> && !is_floating_point_v<U>, int> = 0>
__int128_t dot(Point<T> a, Point<U> b) {
	return (__int128_t) a.x * b.x + (__int128_t) a.y * b.y;
}
template<class T, class U,
	enable_if_t<is_floating_point_v<T> || is_floating_point_v<U>, int> = 0>
auto dot(Point<T> a, Point<U> b) {
	return a.x * b.x + a.y * b.y;
}

// 9) Line / gen_line_from_general / rng: 不依赖 C++17, 照搬正文「点线基础模板」即可
//    (此处省略, 参见 §计算几何/点线基础模板)
	\end{minted}

	\newpage
